% Version 3.5: corrected AS capacity settlement; revalidated March-July and weekly studies.
\documentclass[final,5p,times]{elsarticle}

\usepackage[utf8]{inputenc}
\usepackage{amsmath,amssymb,amsfonts}
\usepackage{booktabs}
\usepackage{graphicx}
\usepackage{xcolor}
\usepackage{url}
\usepackage{array}
\usepackage{tabularx}
\usepackage{algorithm}
\usepackage{algpseudocode}
\usepackage[hidelinks]{hyperref}
\IfFileExists{stfloats.sty}{\usepackage{stfloats}}{}
\renewcommand{\ttdefault}{cmtt}
\setlength{\emergencystretch}{3em}

\journal{Sustainable Energy, Grids and Networks}
\graphicspath{{figures/}}
\biboptions{sort&compress}

\newcommand{\dt}{\Delta t}
\newcommand{\EV}{\mathrm{EV}}
\newcommand{\BESS}{\mathrm{BESS}}
\newcommand{\DA}{\mathrm{DA}}
\newcommand{\RT}{\mathrm{RT}}
\newcommand{\GI}{\mathrm{GI}}
\newcommand{\PV}{\mathrm{PV}}
\newcommand{\BTM}{\mathrm{BTM}}
\newcommand{\TOU}{\mathrm{TOU}}
\newcommand{\PD}{\mathrm{PD}}
\newcommand{\NCD}{\mathrm{NCD}}
\newcommand{\SOC}{\mathrm{SOC}}
\newcommand{\RU}{\mathrm{RU}}
\newcommand{\RD}{\mathrm{RD}}
\newcommand{\SP}{\mathrm{SP}}
\newcommand{\NSP}{\mathrm{NSP}}
\newcommand{\WM}{\mathrm{WM}}
\newcommand{\NWM}{\mathrm{NWM}}
\newcommand{\net}{\mathrm{net}}
\newcommand{\ch}{\mathrm{ch}}
\newcommand{\dch}{\mathrm{dch}}
\newcommand{\actual}{\mathrm{actual}}
\newcommand{\hist}{\mathrm{hist}}
\newcommand{\req}{\mathrm{req}}
\newcommand{\load}{\mathrm{load}}
\newcommand{\calT}{\mathcal{T}}
\newcommand{\calH}{\mathcal{H}}
\newcommand{\calI}{\mathcal{I}}
\newcommand{\calK}{\mathcal{K}}
\newcommand{\pos}[1]{\left[#1\right]_{+}}
\providecommand{\sep}{; }

\begin{document}

\begin{frontmatter}

\title{Retail--Wholesale Value Stacking for Campus Electric-Vehicle Charging, Battery Storage, Photovoltaic Generation, and Building Demand: \\
A Two-Stage 24-h Rolling-Horizon Model Predictive Control Application with Dynamic Forecast Updating}

\author[ucsd]{Yizhan Gu}
\ead{yizhangu@ucsd.edu}
\author[totalenergies]{Wente Zeng}
\ead{wente.zeng@totalenergies.com}
\author[totalenergies]{Hadi Eshraghi}
\ead{hadi.eshraghi@external.totalenergies.com}
\author[ucsd]{Jan Kleissl\corref{cor1}}
\ead{jkleissl@ucsd.edu}
\cortext[cor1]{Corresponding author.}

\affiliation[ucsd]{organization={Center for Energy Research, University of California San Diego},
  city={La Jolla}, postcode={92093}, state={CA}, country={USA}}
\affiliation[totalenergies]{organization={TotalEnergies},
  city={Houston}, state={TX}, country={USA}}

\begin{highlights}
\item EV charging and battery storage are co-optimized for retail--wholesale value stacking.
\item CAISO-inspired DA--RT energy and AS coordination supports portfolio value stacking.
\item Two-stage optimization with real-time  EV forecast update is applied.
\item Five billing months reveal site-dependent value stacking and forecast effects.
\end{highlights}

\begin{abstract}
Efficient use of existing distributed resources is increasingly important as electricity demand and low-carbon generation grow. This paper develops a mixed-integer model predictive control (MPC) framework for a behind-the-meter (BTM) campus portfolio comprising unidirectional electric-vehicle (EV) charging and a bidirectional battery energy storage system (BESS). The portfolio co-optimizes San Diego Gas \& Electric (SDG\&E) retail charges with California Independent System Operator (CAISO) wholesale energy and ancillary-service (AS) revenue. The model follows the chronological day-ahead (DA) and real-time (RT) market process: hourly DA energy and AS awards become fixed commitments, signed 15-min RT recourse modifies the market position, and the resulting obligation is reconciled with EV--BESS capability and the point-of-common-coupling (PCC) meter. A dynamic EV forecast updater converts session arrivals, departures, requested energy, and charger availability into successive 24-h controller inputs. Measured photovoltaic (PV) generation and building demand are introduced separately as passive PCC quantities. Across March--July 2025 at two measured UCSD sites, PV reduces site retail cost and coordinated EV--BESS operation adds further value, while the direction and magnitude of forecast value depend on the site's power scale and carried billing state. The contribution is an application-oriented model that integrates heterogeneous BTM resources, two-stage market chronology, dynamic forecast updating, and executable retail--wholesale value stacking in one auditable framework.
\end{abstract}

\begin{keyword}
Electric vehicle charging \sep Battery energy storage \sep Model predictive control \sep Value stacking \sep Ancillary services \sep Wholesale markets \sep Behind-the-meter resources
\end{keyword}

\end{frontmatter}

\onecolumn
\section*{Nomenclature}
\label{sec:nomenclature}

The following notation is used throughout. Power, energy, state of charge, and
monetary quantities are expressed in kW, kWh, fractions, and US dollars,
respectively, unless a different unit is stated. Superscripts identify a
resource, market stage, service, or accounting channel; subscripts identify an
interval, EV session, AS product, or rolling-controller call.

\begin{table}[H]
\centering
\caption{Abbreviations used in the manuscript.}
\label{tab:abbreviations}
\scriptsize
\begin{tabularx}{\textwidth}{@{}p{0.11\textwidth}X p{0.11\textwidth}X@{}}
\toprule
Abbreviation & Definition & Abbreviation & Definition \\
\midrule
AGC & automatic generation control & AS & ancillary service \\
ASMP & ancillary-service marginal price & BESS & battery energy storage system \\
BTM & behind the meter & CAISO & California Independent System Operator \\
DA & day ahead & DER & distributed energy resource \\
EV & electric vehicle & FERC & Federal Energy Regulatory Commission \\
GI & grid interchange at the PCC & LMP & locational marginal price \\
MILP & mixed-integer linear program & MPC & model predictive control \\
NCD & non-coincident demand charge & NSP & non-spinning reserve \\
NWM & non-wholesale-market channel & PCC & point of common coupling \\
PD & peak-period demand charge & PV & photovoltaic generation \\
RD & regulation down & RT & real time \\
RU & regulation up & SDG\&E & San Diego Gas \& Electric \\
SOC & state of charge & SP & spinning reserve \\
TOU & time of use & UCSD & University of California San Diego \\
V1G & unidirectional managed EV charging & V2G & vehicle to grid \\
VPP & virtual power plant & WM & wholesale-market channel \\
\bottomrule
\end{tabularx}
\end{table}

\begin{table}[H]
\centering
\caption{Principal sets, parameters, states, and decision variables.}
\label{tab:nomenclature}
\scriptsize
\begin{tabularx}{\textwidth}{@{}p{0.20\textwidth}X@{}}
\toprule
Symbol & Definition \\
\midrule
$t\in\calT$, $i\in\calI$, $x\in\calK$ & 15-min interval, EV session, and AS product; $\calK=\{\RU,\RD,\SP,\NSP\}$. \\
$k$, $\calH_k$, $H(t)$ & Rolling-controller call, its active 24-h horizon, and the DA trade hour containing interval $t$. \\
$\mathcal I_k^\pi$, $\omega_k$ & Information available under forecast updater $\pi$ and the realized data observed after interval $k$. \\
$a_i,d_i,E_i^{\req},\bar P_i$ & Arrival, departure, requested energy, and charging-power limit of EV session $i$. \\
$b_{t,i}$, $p_{t,i}^{\EV}$, $p_t^{\EV}$ & EV availability, session charging power, and aggregate nonnegative EV power. \\
$B_t^{\EV}$, $P_{t,\max}^{\EV}$ & Historical EV reference baseline and connected-fleet charging capability. \\
$p_t^{\ch,\BESS}$, $p_t^{\dch,\BESS}$, $s_t$ & BESS charging power, discharging power, and SOC. \\
$C^{\BESS}$, $P^{\BESS}$, $\gamma$ & BESS energy rating, power rating, and round-trip efficiency. \\
$p_t^{\load}$, $p_t^{\PV}$ & Measured passive building demand and PV generation in the sensitivity case. \\
$p_t^{\GI}$ & Signed PCC power; positive denotes retail import and negative denotes export. \\
$p_t^{\DA}$, $p_t^{\RT}$, $p_t^{\actual}$ & Fixed DA energy award, signed RT adjustment, and actual energy position, with $p_t^{\actual}=p_t^{\DA}+p_t^{\RT}$. \\
$c_t^{x,\DA}$, $c_t^{x,\RT}$, $c_t^{x,\actual}$ & Fixed DA AS award, signed RT adjustment, and actual nonnegative capacity for product $x$. \\
$\alpha_t^x$, $q_t^{\WM}$ & AS attenuation factor and the attenuation-adjusted EV--BESS wholesale energy obligation. \\
$\bar P_t^{\mathrm{up}}$, $\bar P_t^{\mathrm{down}}$ & Shared upward and downward capability of the controllable EV--BESS portfolio. \\
$M_k^{q,\hist}$, $z_k^q$, $m_q(t)$ & Carried billing peak, incremental peak variable, and tariff-eligibility indicator for $q\in\{\PD,\NCD\}$. \\
$C_t^{\TOU}$, $C^{\PD}$, $C^{\NCD}$ & Retail energy, peak-period demand, and non-coincident demand rates. \\
$\widetilde\lambda_t^s$, $\widetilde\rho_t^{x,s}$ & Native energy LMP and ASMP for stage $s\in\{\DA,\RT\}$. \\
$\lambda_t^s$, $\rho_t^{x,s}$ & Price coefficients normalized to the common US\$/kWh-equivalent objective basis. \\
$R^{\EV}$, $R^{\WM}$, $V$, $J$ & EV charging-service revenue, wholesale-market revenue, net portfolio value, and minimized cost, with $V=-J$ for like-for-like comparisons. \\
\bottomrule
\end{tabularx}
\end{table}

\twocolumn

\section{Introduction}
\label{sec:introduction}

Electricity systems are being transformed simultaneously by transportation electrification, rapid growth in digital and artificial-intelligence loads, and the policy-driven expansion of low-carbon generation. Meeting this growth only by adding supply and network capacity is costly and slow. A complementary strategy is to operate existing electricity infrastructure more efficiently: flexible demand, storage, and local generation can reshape net load, absorb renewable production, reduce coincident stress, and supply operating reserves. Such coordination supports decarbonization while improving reliability and limiting avoidable capacity expansion. This system-level motivation is reflected in US grid-modernization initiatives and in Federal Energy Regulatory Commission (FERC) Order No.~2222, which seeks to remove barriers preventing aggregations of distributed energy resources (DERs) from participating in organized wholesale markets \cite{doeGridModernization} \cite{ferc2222}.

Behind-the-meter (BTM) portfolios are a natural place to realize this flexibility. Commercial and residential sites increasingly contain electric-vehicle (EV) charging, battery energy storage systems (BESS), photovoltaic (PV) generation, and large building loads. These resources share an interconnection, but they have different physical roles. A BESS can absorb or deliver power subject to state of charge (SOC), efficiency, and cycling limits. Workplace EVs without vehicle-to-grid (V2G) hardware cannot export battery energy, yet their charging can be advanced or delayed while vehicles are plugged in. PV and ordinary building demand are not dispatchable at the timescale considered here, but they change the net power measured at the point of common coupling (PCC). Evaluating each resource in isolation can therefore overstate value: a profitable wholesale action may simultaneously establish a higher retail demand-charge peak at the shared meter.

The economic opportunity is \emph{value stacking}: one coordinated portfolio reduces retail electricity cost while participating in day-ahead (DA) and real-time (RT) wholesale energy and ancillary-service (AS) markets. The retail cost consists of energy and demand charges: retail time-of-use (TOU) prices reward shifting import away from expensive periods, whereas monthly peak-demand (PD) and non-coincident-demand (NCD) charges depend on the largest metered import reached over a billing period. Wholesale prices can instead reward injection, absorption, or the reservation of upward and downward capability, and the signs and timing of these signals can conflict. Consequently, the relevant quantity is the portfolio's net value after all charges and settlements.

Market chronology adds a further constraint. The California Independent System Operator (CAISO) procures regulation up (RU), regulation down (RD), spinning reserve (SP), and non-spinning reserve (NSP) in addition to energy \cite{caisoAncillary2026}. In the modeled two-stage sequence, cleared DA awards are binding commitments and signed RT decisions increase or decrease the corresponding positions; actual delivery is their sum. This abstraction is consistent with CAISO documentation that distinguishes DA and RT AS awards and reserves physical capability for financially binding DA awards \cite{caisoMarketOperations2026} \cite{caisoNGRCommitment2020}. The resulting obligation must fit the common capability of the connected EV fleet and BESS. Deployment would additionally require scheduling-coordinator representation, resource certification, telemetry, metering, and distribution-utility coordination \cite{caisoDER2026}. The model represents the operational chronology and financial coupling; it does not replace those eligibility and performance requirements.

Forecast uncertainty is inseparable from this chronology. Workplace charging availability and requested energy depend on stochastic arrivals, departures, and user behavior. A useful controller must act before all sessions are known, preserve commitments that have already cleared, and update the remaining schedule as measurements arrive. MPC provides this repeated decision structure. At each 15-min interval, the controller constructs a 24-h horizon, calls the optimizer, implements only the current action, updates physical and billing states, and rolls the horizon forward. The term \emph{rolling horizon} is used throughout; it is the same control principle commonly called a receding horizon. Two EV forecast updaters are compared inside this controller. Persistence reuses recent observed session patterns and is deployable. Perfect information supplies realized future EV sessions only within the same 24-h window and isolates EV forecast error; it is not a claim of practical predictability. A separate full-horizon benchmark knows the complete selected evaluation week. It is an ex-post information benchmark; an upper-bound interpretation additionally requires a common feasible set, including the DA offer envelope.

The study is built around measured UC San Diego (UCSD) workplace-charging data and a co-located BESS. EVs and the BESS are the controllable market portfolio. Measured PV generation and building demand are introduced in a separate sensitivity as passive additions to the signed retail meter. They affect retail cost and can indirectly change optimal EV--BESS actions, but they do not contribute wholesale bids because no independent controllability or guaranteed response is assumed.

The main contribution of this work is an integrated framework consisting of these closely related components:
\begin{enumerate}
  \item \textbf{Integrated retail--wholesale value stacking.} A single mixed-integer formulation co-optimizes unidirectional EV flexibility and bidirectional BESS operation across EV service, retail TOU and monthly demand charges, wholesale energy, and four AS products. A shared deliverability envelope prevents the same physical capability from being allocated to incompatible services.
  \item \textbf{Market-sequence-consistent control.} The model follows the chronological transition from DA submission to RT recourse. Cleared DA energy and AS awards are fixed, signed RT quantities represent market increases or decreases, and actual nonnegative AS capacity and deployed energy are reconciled with physical power.
  \item \textbf{Forecast-integrated 24-h MPC framework.} A dynamic forecast-to-control layer translates updated EV session information into availability, requested-energy, and charging-capability inputs at every rolling optimization. The deployable Persistence and Perfect-information updaters are implemented through an otherwise identical MPC workflow, so their executed dispatch and economic outcomes measure the operational effect of the forecast update.
  \item \textbf{Site- and season-resolved application evidence.} Five consecutive billing months at two measured sites separate the value of passive PV from the full EV--BESS portfolio. Monthly financial decomposition and daily dispatch reveal how site power scale, seasonal tariffs, and carried billing states change the value of the EV forecast updater.
  
\end{enumerate}

The remainder of the paper follows the decision process of an operator. Section~\ref{sec:related_work} reviews EV aggregation, multi-resource portfolios, value stacking, market participation, and MPC. Sections~\ref{sec:system}--\ref{sec:mpc} describe the physical system, data pipeline, forecast updating, and market timeline. Section~\ref{sec:model} presents the compact optimization, while \ref{app:complete_model} records the complete formulation. Sections~\ref{sec:case_study} and \ref{sec:accounting} define the case study, financial reconstruction, and validation protocol. Section~\ref{sec:results} reports monthly and daily results, and the final sections discuss interpretation, limitations, and implications for larger BTM portfolios.

\section{Related work and research gap}
\label{sec:literature}
\label{sec:related_work}

\subsection{Unidirectional EV charging flexibility and aggregation}
Coordinated EV charging has been studied for distribution-impact mitigation, valley filling, renewable integration, and charging-cost reduction \cite{clement2010impact} \cite{sortomme2011optimal} \cite{gan2013optimal} \cite{richardson2013electric}. Its feasible set changes when vehicles arrive and depart, and each session has an individual energy requirement. Unidirectional charging has been shown to provide scheduled flexibility without exporting vehicle energy \cite{sortomme2011optimal}. Compact aggregation models have subsequently represented heterogeneous EV populations \cite{zheng2013aggregation} \cite{nazari2022updated} \cite{alTaha2025aggregateFlex}, while probabilistic formulations quantify time and energy flexibility before departure rather than reducing uncertainty to one expected load trace \cite{huber2020probabilistic}.

For a unidirectional fleet, upward grid flexibility is a reduction of charging relative to a feasible reference, and downward flexibility is an increase in charging. Market-facing EV research has optimized aggregator participation in energy and reserve markets under uncertain availability and deviations \cite{vagropoulos2013optimal} \cite{bessa2014optimization}. Practical demonstrations have also shown the potential of managed EV fleets to provide AS \cite{deforest2018lafb}. Two recent UCSD studies are the closest methodological predecessors. Workplace charging flexibility was first optimized against TOU energy, demand charges, and a DA demand-response opportunity \cite{chen2024costoptimal}. The subsequent framework added a two-layer DA/RT architecture in which online MPC updates charging timing and user-service flexibility to mitigate EV forecast errors \cite{chen2026realtime}. That predecessor, however, considers EV charging in a demand-response setting without a co-optimized BESS, a retail--wholesale energy ledger, or four CAISO AS products. The present work retains its service-preserving, measurement-updated EV structure and extends the market boundary to an EV--BESS portfolio with fixed DA awards, signed RT recourse, and one PCC retail account.

\subsection{EV--BESS portfolios and virtual-storage interpretations}
Combining EV charging with stationary storage is attractive because the resources are complementary. The BESS supplies bidirectional, fast power but has a finite energy state and cycling cost; EV charging supplies a changing envelope of controllable demand. Earlier work has optimized EV charging with energy storage under electricity-market prices \cite{jin2013evstorage} and has assessed EV fleets as aggregated storage and flexibility resources \cite{mills2018evstorage}.

A broader virtual power plant (VPP) literature considers aggregation of heterogeneous DERs for coordinated market participation \cite{pudjianto2007vpp} \cite{zamani2016vpp} \cite{qiu2017transactive}. The VPP concept provides an institutional and systems perspective: distributed resources can be represented to the market through an aggregator while their internal constraints remain coordinated. Existing formulations span several distinct objectives. Some emphasize system integration and the aggregation of DER capabilities for system-level operation \cite{pudjianto2007vpp}; others formulate stochastic scheduling of large-scale VPPs participating in energy and reserve markets \cite{zamani2016vpp}; and transactive-energy formulations coordinate DER operation across wholesale and distribution-level settings \cite{qiu2017transactive}. In contrast, these formulations do not explicitly reconstruct commercial retail demand charges from the metered load resulting from wholesale-market actions. The present study considers a physically behind-the-meter campus portfolio in which wholesale decisions directly alter the customer's metered demand and therefore its retail bill, requiring wholesale-market participation and retail demand charges to be modeled jointly.

\subsection{BTM tariffs and multi-service value stacking}
BTM storage is commonly scheduled for energy arbitrage, peak shaving, renewable self-consumption, and grid services \cite{raoufat2018bess} \cite{wu2017analytical} \cite{chen2021valueStacking}. Value stacking increases potential asset utilization, but the individual services compete for the same power and stored energy. Revenue stacking across BTM energy and AS markets has been studied explicitly \cite{seward2022revenueStacking}, while distributed control and storage-sizing studies show how tariff and service design determine which combinations are economically meaningful \cite{karandeh2022distributed} \cite{zhang2024btmsizing}.

Retail demand charges make this coupling path dependent. The charge is based on a billing-period (monthly) maximum, so an action taken early in the month can establish a threshold that changes the marginal cost of all later actions. Behind-the-meter PV and storage economics are likewise sensitive to the retail-rate structure \cite{boampong2020retail}. These findings motivate the incremental demand-charge state used in this paper. They also explain why a 24-h controller cannot generally reproduce the full billing-period optimum even when it knows the EV sessions inside its local horizon.

\subsection{Joint retail--wholesale participation and market fidelity}
Microgrid energy-management research has co-optimized local resources under market uncertainty and operational risk \cite{parisio2014model} \cite{shen2016microgrid} \cite{hasankhani2021stochastic}. Other studies have coordinated renewable generation, BESS, energy trading, and ancillary services \cite{santosuosso2024sempc}, or have combined PV, flexible loads, and storage for frequency reserves \cite{warmerdam2026power}. These contributions establish the value of multi-resource scheduling, but the settlement boundary matters. A customer connected behind a retail meter cannot evaluate wholesale bids independently of TOU and demand-charge cost, and passive PV or building demand cannot be credited as callable capacity without a control and verification mechanism.

CAISO's market design adds a specific chronological and institutional structure. Energy and AS positions pass from DA awards to RT adjustments, while actual AS capacity must remain nonnegative and physically deliverable. The four modeled AS products have different directions and response characteristics \cite{caisoAncillary2026}. DER aggregations may participate in CAISO energy and AS markets, but actual deployment requires certified resources, scheduling-coordinator representation, metering, telemetry, and distribution-utility coordination \cite{caisoDER2026}. Many optimization papers approximate one or more of these elements. The contribution here is not a claim of complete settlement replication; it is a transparent operational abstraction that keeps the irreversible DA/RT sequence, signed recourse, deployed-energy accounting, and common EV--BESS capability in one auditable model.

\subsection{Forecast updating and rolling-horizon control}
MPC is well suited to microgrids and charging stations because it repeatedly incorporates new information while respecting dynamic constraints \cite{parisio2014model} \cite{rawlings2017mpc} \cite{hermans2024implementation} \cite{klemets2024charging}. PV--EV--BESS studies have demonstrated explicit and robust MPC under uncertain renewable and charging conditions \cite{cabrera2022realtime} \cite{cabrera2024robust}. These works establish MPC as a mature control framework; the present novelty is how the information update is embedded in a retail--wholesale commitment and settlement model.

Forecast quality and controller value are not interchangeable. A forecast with lower statistical error does not necessarily produce higher closed-loop economic value because the objective is asymmetric, constrained, and state dependent \cite{ludolfinger2025prediction}. Moreover, Perfect information inside a truncated horizon is not a full-horizon optimum. If the controller omits the continuation value of SOC, demand thresholds, unfinished EV service, and committed bids beyond its local view, an attractive first action can reduce later value. Three concepts are therefore kept distinct: a deployable Persistence updater, Perfect information within the same 24-h controller, and a full-horizon benchmark whose upper-bound interpretation is conditional on common feasible-set assumptions.

\subsection{Research gap and paper novelty}
Table~\ref{tab:literature_gap} positions the work against representative application classes. Existing literature treats EV aggregation, BTM storage, VPP coordination, retail tariffs, wholesale energy and reserves, and MPC with considerable depth. Much less attention has been paid to their intersection: a unidirectional workplace EV fleet and BESS behind one demand-charged meter; fixed DA awards followed by signed RT recourse; non-double-counted energy and AS deliverability; forecast updates evaluated from executed decisions; and passive PV/building profiles that alter only the retail meter. That intersection is the paper's novelty boundary.

\begin{table*}[!t]
\centering
\caption{Representative literature streams and the unresolved coupling addressed in this paper.
``Retail state'' denotes explicit TOU energy charges and billing-period demand-charge accounting.
``DA/RT + AS'' denotes joint representation of day-ahead commitments, real-time recourse or settlement, and ancillary services.
``Executed MPC'' denotes rolling-horizon control updated using realized system information.}
\label{tab:literature_gap}
\scriptsize
\setlength{\tabcolsep}{3.0pt}

\begin{tabularx}{\textwidth}{
@{}
>{\raggedright\arraybackslash}p{0.205\textwidth}
>{\centering\arraybackslash}p{0.085\textwidth}
>{\centering\arraybackslash}p{0.105\textwidth}
>{\centering\arraybackslash}p{0.125\textwidth}
>{\centering\arraybackslash}p{0.105\textwidth}
>{\raggedright\arraybackslash}X
@{}}
\toprule

Study class
& Uni-dir. EV
& Retail state
& DA/RT + AS
& Executed MPC
& Key contrast \\

\midrule

EV aggregation and reserve bidding
\cite{vagropoulos2013optimal} \cite{bessa2014optimization}
&
$\checkmark$
&
--
&
DA energy +\newline reserve/recourse
&
RT control,\newline not MPC
&
No retail demand-charge state or BESS coupling. \\

UCSD workplace charging
\cite{chen2024costoptimal} \cite{chen2026realtime}
&
$\checkmark$
&
$\checkmark$
&
--
&
$\checkmark$
&
No combined DA/RT energy and AS stack. \\

BTM BESS market stacking
\cite{hanif2022multiservice} \cite{seward2022revenueStacking}
&
--
&
Demand/energy\newline charges
&
Energy + AS
&
--
&
Stationary-storage service stacking without shared EV flexibility or executed retail-aware MPC. \\

BTM BESS retail MPC
\cite{idomar2026stacked}
&
--
&
$\checkmark$
&
aFRR +\newline RT operation
&
$\checkmark$
&
Retail-aware BESS MPC without unidirectional EV flexibility or fixed DA AS commitments. \\

PV--EV--BESS MPC
\cite{cabrera2022realtime} \cite{cabrera2024robust}
&
$\checkmark$
&
Energy-price\newline signals only
&
--
&
$\checkmark$
&
No billing-period demand charges or wholesale market stack. \\

VPP system integration
\cite{pudjianto2007vpp}
&
--
&
--
&
--
&
--
&
System-level DER aggregation rather than coupled retail/wholesale scheduling. \\

Stochastic VPP scheduling
\cite{zamani2016vpp}
&
--
&
DR only
&
DA energy +\newline reserve
&
--
&
Day-ahead heterogeneous-DER scheduling without RT recourse or retail demand charges. \\

Transactive VPP scheduling
\cite{qiu2017transactive}
&
--
&
Retail\newline energy rate
&
DA/RT energy;\newline no AS
&
RT rescheduling,\newline not MPC
&
DA/RT heterogeneous-DER coordination without AS or commercial demand-charge accounting. \\

Renewable--storage market MPC
\cite{santosuosso2024sempc}
&
--
&
--
&
Energy + AS
&
$\checkmark$
&
Renewable generation and storage trading without EV or retail demand-charge coupling. \\

PV/load/BESS reserve control
\cite{warmerdam2026power}
&
V2G
&
--
&
AS provision
&
$\checkmark$
&
Reserve control with bidirectional EV flexibility rather than unidirectional workplace charging. \\

This work
&
$\checkmark$
&
$\checkmark$
&
$\checkmark$
&
$\checkmark$
&
Joint retail state, DA/RT recourse, AS, and shared EV--BESS capability. \\

\bottomrule
\end{tabularx}
\end{table*}

The claimed contribution is an application-level integration: a measured workplace portfolio, a retail tariff with carried billing states, a chronological CAISO-inspired market layer, and a closed-loop forecast comparison whose financial statements can be reconstructed from executed 15-min trajectories.

\section{System description and data flow}
\label{sec:system}

\subsection{Campus BTM portfolio}
Figure~\ref{fig:architecture} separates the physical PCC from
the market-aggregation function. The PCC is the metered
electrical boundary at which $p_t^{\GI}$ is measured and the
retail tariff is applied. The EV--BESS market aggregator is
instead a logical control and scheduling interface,
implemented through the rolling 24-h MPC; it is analogous
to a small VPP. The controllable portfolio consists of a
workplace EV aggregation and a BESS behind the common meter.
Although power is aggregated at each modeled site boundary,
each EV session retains its own arrival, departure, requested
energy, and interval charging limit. The BESS is represented
by charge/discharge power, SOC dynamics, round-trip
efficiency, throughput limits, and a terminal-state condition.

The information and control connections are distinguished
from the physical power flows. The EV fleet supplies
session information, charger limits, and inputs from the
Persistence or Perfect forecast updater, while the BESS
supplies its SOC and operating limits. The MPC returns
charging commands to the EV fleet and charge/discharge
commands to the BESS. Retail tariffs, wholesale prices,
AS attenuation factors, and cleared market awards provide
additional optimization inputs. Executed charging, battery
SOC, metered demand, and existing commitments are carried
forward to update subsequent controller decisions.

PV generation and building demand are measured, exogenous
profiles in the passive-resource comparisons. Because the
controller cannot curtail PV, shift building demand, or
guarantee their response, they are not assigned energy bids,
AS-capacity offers, or wholesale revenue. Their profiles
are supplied to the MPC to account for their contribution
to net metered demand and retail cost, but no dispatch
commands are returned to these resources. This modeling
choice asks a focused sensitivity question: how does the
EV--BESS controller behave when the retail net-load
environment changes?

\begin{figure*}[!t]
  \centering
  \includegraphics[width=0.96\textwidth]{fig_01_system_architecture_v31.png}
  \caption{Physical, control, and market interfaces of the
  campus BTM portfolio. EV and BESS information enters the
  rolling 24-h MPC, which returns resource dispatch commands
  and determines DA bids and RT adjustments for CAISO energy
  and AS participation. Building demand and PV generation
  are passive inputs without wholesale offers. All four
  resources share the electrical connection to the PCC,
  where net grid interchange determines retail settlement.
  Dashed arrows denote information or settlement data,
  colored solid arrows denote decisions, and dark solid
  arrows denote physical power flows.}
  \label{fig:architecture}
\end{figure*}

The portfolio has three coupled boundaries. First, the
\emph{service boundary} requires every included EV to
receive its requested energy before departure. Second,
the \emph{physical boundary} limits total upward and
downward response by the simultaneous EV and BESS
capability. Third, the \emph{financial boundary} reconciles
wholesale actions with the same signed meter used for
retail TOU and demand charges. Together, these boundaries
require market schedules to respect the modeled
charging-service and device constraints while accounting
for their associated retail costs.

\subsection{Retail and wholesale interfaces}
The retail account uses signed interval energy at the PCC. Import has positive cost, while export is credited at the same modeled TOU coefficient; this symmetric export treatment is a study assumption rather than a universal SDG\&E rule. The modeled large-commercial tariff distinguishes summer and winter TOU periods, applies peak-period demand charges only within eligible windows, and assesses NCD from the maximum billing-cycle demand \cite{sdgeLargeCommercial2025} \cite{sdgeDemandGuide2026}. Both demand charges are stateful over the monthly billing period.

The wholesale interface contains DA and RT energy and four AS products. RU, SP, and NSP reserve upward capability; RD reserves downward capability. CAISO describes regulation as automatic response for frequency control and requires participating resources to be certified and capable of delivering the awarded service; SP and NSP must be available within the specified reserve response time \cite{caisoAncillary2026}. The energy position and AS deployment therefore share one headroom envelope. For example, the portfolio cannot schedule all upward power as energy and sell the same capability again as RU. Capacity payment compensates reserved capability, while expected deployed energy is valued at the RT energy price. The model consequently distinguishes an AS capacity award from the energy obligation created when a fraction of that award is activated.

\subsection{Data sources and temporal alignment}
All operating quantities are aligned to a common 15-min index,
with $\dt=0.25$~h. The EV representation builds on the
session-based framework in \cite{chen2026realtime}, in which
arrival time, plug-in duration, and session energy characterize
individual charging requirements on a discrete-time grid.
These session-level quantities are retained in the present
study, while site-specific selection and an explicit
energy-feasibility reconciliation are used to prepare the
inputs for the integrated EV--BESS optimization.

The raw EV data comprise session reports and interval
measurements. Session reports provide connection and
disconnection timestamps, delivered energy, and user and
station identifiers, whereas interval records provide measured
energy and average and maximum charging power. Duplicate
records are removed from each report, timestamps are converted
to a common datetime representation, and interval measurements
are linked to their corresponding sessions using the unique
session identifier. Records without the required timestamps,
identifiers, or electrical measurements are excluded. The
retained dataset is restricted to the modeled charger types
and same-day sessions; overnight sessions are not included in
the present implementation. Sessions with less than
1~kWh of delivered energy or less than 15~min of charging
duration are also excluded.

Measurement consistency is checked before the optimization
inputs are assembled. For a 15-min interval, measured energy
in kWh divided by $0.25$~h is compared with the reported
average power in kW, and records with a relative discrepancy
of 10\% or greater are excluded. Session energy is reconciled
with the sum of its processed interval-energy measurements.
For the site-specific experiments, station metadata are then
used to select the EVSEs belonging to each study portfolio;
the resulting charger counts are checked against the selected
site configuration. This prevents charging sessions from
unrelated campus locations from entering the site-level
optimization.

A further consistency check is applied when the session data
are mapped to the optimization grid. For each session, the
maximum deliverable energy is obtained by summing its
interval-specific charging-energy upper bounds over the
represented connection period. If the processed session-energy
target exceeds this envelope, the target is capped at the
maximum deliverable energy before forecasts are constructed.
This reconciliation prevents inconsistent measurement or
time-discretization inputs from making the charging-service
equality incompatible with the interval charging limits.
It is an input-data correction, rather than an additional
decision allowing the optimizer to curtail charging demand.

Wholesale inputs consist of CAISO hourly DA and 15-min RT energy prices and 1-h DA/15-min RT AS capacity prices aligned to the same 15-min time step. Hourly DA awards are
repeated and constrained equal within each four-interval block, whereas dispatchable RT AS awards remain interval-specific, consistent with CAISO's published market-instrument definitions \cite{caisoMarketInstruments2026} \cite{caisoMarketOperations2026}.
Native market prices are converted to the common US\$/kWh-equivalent coefficients derived in Section~\ref{sec:price_resolution} before entering the
objective or a figure. Retail prices use the implemented summer/winter TOU schedule and PD/NCD rates; their period definitions and demand-charge interpretation follow the published SDG\&E large-commercial materials \cite{sdgeLargeCommercial2025} \cite{sdgeDemandGuide2026}.

Measured PV generation and building-meter power are aligned
to the common 15-min grid. The source timestamps denote
interval ends and are shifted backward by 15~min to represent
interval starts consistently with the optimization inputs.
Duplicate timestamps are rejected, and missing measurements
are identified after reindexing. Internal gaps are interpolated
in time with a limit of eight consecutive 15-min entries;
remaining missing values are flagged and must be resolved
before the data are used. Negative PV measurements are
clipped to zero, whereas signed building-meter measurements
are preserved because negative net load can represent
electricity export.

The building measurements represent net demand after the
contribution of the PV systems included within their metering
boundaries. Gross building demand is therefore reconstructed
by adding the corresponding measured PV generation to the
net-meter readings. Only PV systems mapped to those building
meters are included in this reconstruction. In the subsequent
PCC power balance, gross building demand is added and PV
generation is subtracted once, thereby recovering the
measured building--PV net demand without double-counting
PV production. The selected building and PV series are
aggregated separately for each site and remain exogenous
inputs to the optimization.

Regulation-attenuation factors are processed separately from
official CAISO hourly attenuation tables and repeated across
the four constituent 15-min intervals. These factors represent
expected deployment proxies rather than realized
automatic-generation-control (AGC) trajectories.

\subsection{Data flow and reproducibility boundary}
For each evaluation day, the pipeline constructs realized EV sessions, forecast sessions, carried BESS and billing states, market prices, and any active passive-meter profiles. The DA optimizer produces commitments. Each RT optimization call reads those commitments as constants, receives the updated EV information, and implements one interval. Only the first executed quantities update the physical state and financial ledger; the remaining planned entries are discarded when the horizon advances.

\section{EV forecast updater}
\label{sec:forecasting}

\subsection{Forecast objects}
The EV forecast determines the expected sessions, their connection windows, requested energy, and interval charging capability. Let $\widehat{\mathcal I}_{k}^{\pi}$ denote the sessions visible at controller call $k$ under updater $\pi$. These sessions generate forecast availability $\widehat b_{t,i}$, requested energy $\widehat E_i^{\req}$, and aggregate capability $\widehat P_{t,\max}^{\EV}$. Once a vehicle arrives, its measured state replaces the forecast representation in RT. Forecast error therefore changes the remaining feasible set.

\subsection{Persistence updater}
Persistence is the practical low-complexity updater. Before the operating day, it uses the most recent eligible historical day in the corresponding DA/RT construction to form the expected session population and energy requests. During RT execution, observed arrivals, delivered energy, and connected-vehicle limits replace the corresponding forecast entries. If the historical template contains more or fewer vehicles than the realized fleet, session records are trimmed or repeated by charger class to preserve a feasible input dimension. Persistence is intended as a transparent benchmark for future learned forecasts rather than as a statistically optimal predictor.

The EV market baseline averages the preceding ten eligible weekdays or four eligible weekend/holiday days of the same calendar class. Historical intervals already altered by an activated market response are excluded so the reference is not contaminated by earlier control actions. No day-of adjustment is applied in the present study. Thus the baseline is a reference charging profile, whereas Persistence predicts future arrivals and requested energy.

\subsection{Perfect-information updater}
Perfect information supplies the realized sessions and requested energy within the same 24-h look-ahead. It removes EV forecast error but leaves every other controller feature unchanged: horizon length, DA/RT chronology, tariffs, market prices, BESS model, initial state, terminal conditions, and optimizer tolerances are identical. It is therefore an information benchmark for the 24-h controller, not a deployable forecast and not a billing-period upper bound.

This distinction matters because the closed-loop controller is path dependent. A first action changes SOC, the remaining energy of connected EVs, the monthly demand thresholds, and future market flexibility. Better information inside a truncated horizon can improve the local optimization yet still leave a state that is less valuable after the horizon boundary. A full-period oracle removes that boundary. Its exact optimum bounds a controller trajectory only if that trajectory belongs to the same feasible set, with common initial and terminal conditions; the reported weekly benchmark has an additional DA-envelope qualification.

\section{24-h MPC and market execution}
\label{sec:mpc}

\subsection{Day-ahead and real-time chronology}
Figure~\ref{fig:timeline} summarizes the chronology for operating day $d$ (day~0). Historical EV measurements available by day~$-2$ are first assembled for the Persistence updater. On day~$-1$, the updater produces the EV input and the DA optimizer determines hourly energy and AS offers before bid submission. CAISO's published timeline closes DA bids at 10:00 and publishes market results at approximately 13:00; the intermediate preparation times shown in the figure are the study's internal workflow checkpoints \cite{caisoTimeline2026}. Cleared hourly awards are then stored as constants for day~0. DA and RT AS are awarded and RT AS procurement is dispatched on a 15-min basis \cite{caisoMarketOperations2026} \cite{caisoMarketInstruments2026}.

At midnight on day~0, the RT rolling controller begins. At each 15-min interval it updates observed EV sessions, delivered EV energy, BESS SOC, and the billing peaks; calls the 24-h optimizer with the cleared DA quantities fixed; implements the first decision; and advances by one interval. Hence the model respects two distinct clocks: a one-hour DA award that is constant over four model intervals and a 15-min RT recourse decision. The simulation abstracts from CAISO's 5-min dispatch layer and from telemetry/performance settlement, which remain deployment requirements rather than optimization variables in this study.

\begin{figure*}[!t]
  \centering
  \includegraphics[width=0.98\textwidth]{fig_02_legacy_market_bidding_mpc_timeline.png}
  \caption{Market bidding and rolling-horizon execution. DA energy and AS bids are submitted before the operating day and become fixed commitments. RT optimization uses updated EV information to choose signed adjustments at 15-min resolution. Internal preparation times precede the CAISO 10:00 DA gate closure and 13:00 result publication.}
  \label{fig:timeline}
\end{figure*}

The RT variables are defined as changes, not as the complete physical bid. For energy, $p_t^{\actual}=\bar p_t^{\DA}+p_t^{\RT}$, where $\bar p_t^{\DA}$ is the cleared commitment. For every AS product $x$, $c_t^{x,\actual}=\bar c_t^{x,\DA}+c_t^{x,\RT}$. The signed adjustment may reduce a commitment when updated EV availability or BESS state makes the original award less attractive or infeasible, but the resulting actual capacity must remain nonnegative. The objective settles that adjustment economically; it does not use a large artificial penalty to force RT decisions toward DA values.

\subsection{Rolling-horizon state update}
Let the state before call $k$ be
\begin{align}
\xi_k=\left(s_k,\boldsymbol e_k,M_k^{\PD,\hist},M_k^{\NCD,\hist},
\bar p_k^{\DA},\bar{\boldsymbol c}_k^{\DA}\right),
\label{eq:state_vector_main}
\end{align}
where $s_k$ is BESS SOC, $\boldsymbol e_k$ stores delivered EV energy,
$M_k^{\PD,\hist}$ and $M_k^{\NCD,\hist}$ are the carried monthly demand
thresholds, and $\bar p_k^{\DA}$ and $\bar{\boldsymbol c}_k^{\DA}$ are the
active DA energy and AS commitments. Given information $\mathcal I_k^\pi$
from updater $\pi$, the controller computes
\begin{align}
\boldsymbol x_k^\star\in\arg\min_{\boldsymbol x\in\mathcal X(\xi_k,\mathcal I_k^\pi)}
J_k^{24\mathrm h}(\boldsymbol x),
\label{eq:mpc_problem_main}
\end{align}
implements only $x_{k|k}^\star$, and updates
\begin{align}
\xi_{k+1}=F_k(\xi_k,x_{k|k}^\star,\omega_k).
\label{eq:mpc_update_main}
\end{align}
Here $\boldsymbol x_k$ contains the horizon decisions, $\mathcal X$ is the
feasible set defined in Section~\ref{sec:model}, and $J_k^{24\mathrm h}$ is the
active-horizon economic objective in \eqref{eq:objective}. Only the first
decision $x_{k|k}^\star$ is implemented. The transition map $F_k$ then uses
the realized interval data $\omega_k$ to propagate BESS SOC, delivered EV
energy, billing peaks, and active DA commitments. Forecasting creates a forward
plan; measurement closes the loop before the next action.

\subsection{Execution algorithm}
Algorithm~\ref{alg:rolling_mpc} states the complete operating loop. The DA
optimizer and every RT call use the same physical and economic formulation;
the difference is that RT receives fixed DA awards and signed recourse
variables.

\begin{algorithm}[!t]
\caption{Two-stage DA/RT rolling-horizon controller}
\label{alg:rolling_mpc}
\begin{algorithmic}[1]
\Require Evaluation set $\calT$, updater $\pi$, prices and tariff, initial state $\xi_0$
\Ensure Executed dispatch and reconciled financial ledger
\For{each operating day $d$}
  \State construct the DA EV forecast and 24-h physical capability
  \State call the DA optimizer and store hourly awards $\bar p^{\DA}$ and $\bar{\boldsymbol c}^{\DA}$
  \For{each 15-min RT interval $k$ in day $d$}
    \State merge arrivals and measurements available at $k$ with updater $\pi$
    \If{$k$ is the final executed interval of $\calT$}
      \State include the common terminal-SOC constraint
    \EndIf
    \State call the RT optimizer over $\calH_k$ with DA awards fixed
    \State implement $p_{k}^{\EV}$, $p_{k}^{\BESS}$, $p_{k}^{\RT}$, and $\boldsymbol c_{k}^{\RT}$ only
    \State update $\xi_{k+1}$ and the ledger using realized data $\omega_k$
  \EndFor
\EndFor
\State verify terminal SOC and reconcile the final financial ledger
\end{algorithmic}
\end{algorithm}


\section{Optimization model}
\label{sec:model}

\subsection{Sets, variables, and sign conventions}
Let $t\in\calT$ index 15-min intervals, $i\in\calI$ index EV sessions, and
$x\in\calK=\{\RU,\RD,\SP,\NSP\}$ index AS products. Positive
$p_t^{\GI}$ denotes site import and negative $p_t^{\GI}$ denotes export. Positive
market energy $p_t^{\actual}$ denotes upward delivery, while negative values
denote absorption. The positive-part operator is $\pos{z}=\max(z,0)$.

\begin{table*}[!t]
\centering
\caption{Core variables and parameters.}
\label{tab:notation}
\scriptsize
\begin{tabularx}{\textwidth}{@{}p{0.20\textwidth}X@{}}
\toprule
Symbol & Definition \\
\midrule
$p_{t,i}^{\EV}$, $p_t^{\EV}$ & EV-session charging power and fleet total; both are nonnegative.\\
$B_t^{\EV}$, $P_{t,\max}^{\EV}$ & EV baseline power and connected-fleet charging capability.\\
$p_t^{\ch,\BESS}$, $p_t^{\dch,\BESS}$, $s_t$ & BESS charge, discharge, and SOC.\\
$p_t^{\DA}$, $p_t^{\RT}$, $p_t^{\actual}$ & fixed DA energy, signed RT adjustment, and actual energy position.\\
$c_t^{x,\DA}$, $c_t^{x,\RT}$, $c_t^{x,\actual}$ & fixed DA AS, signed RT adjustment, and actual nonnegative AS capacity.\\
$M_k^{q,\hist}$ & executed billing-period threshold entering MPC call $k$, $q\in\{\PD,\NCD\}$.\\
$C_t^{\TOU}$, $C^q$ & retail energy and demand-charge rates.\\
$\lambda_t^s$, $\rho_t^{x,s}$ & Normalized DA/RT energy and AS price coefficients (US\$/kWh equivalent).\\
\bottomrule
\end{tabularx}
\end{table*}

\subsection{EV service and BESS state}
The main text retains the equations that couple user service and storage state
to the market problem. For availability $b_{t,i}\in\{0,1\}$, charger rating
$\bar P_i$, charging efficiency $\eta_i$, requested session energy
$E_i^{\req}$, and connection set $\calT_i=[a_i,d_i)$,
\begin{align}
0&\le p_{t,i}^{\EV}\le \bar P_i b_{t,i},
&\sum_{t\in\calT_i}\eta_i p_{t,i}^{\EV}\dt&=E_i^{\req},
&p_t^{\EV}&=\sum_i p_{t,i}^{\EV}. \label{eq:ev_core}
\end{align}
The first relation limits each session power, the second guarantees its
requested energy by departure, and the third aggregates the campus fleet. Thus
EV ``discharge'' in the accounting layer is only a reduction from the
nonnegative baseline $B_t^{\EV}$; no vehicle exports power. For BESS energy
rating $C^{\BESS}$ and round-trip efficiency $\gamma$, the SOC evolves as
\begin{align}
s_{t+1}=s_t+\frac{\dt}{C^{\BESS}}
\left(\sqrt{\gamma}\,p_t^{\ch,\BESS}
-\frac{p_t^{\dch,\BESS}}{\sqrt{\gamma}}\right), \label{eq:soc}
\end{align}
where $p_t^{\ch,\BESS}$ and $p_t^{\dch,\BESS}$ are nonnegative charge and
discharge powers. SOC bounds, direction exclusivity, power limits,
daily-throughput, and terminal conditions are given in \ref{app:complete_model}.

\subsection{Meter balance and passive DER extension}
\label{sec:passive_der}
The general signed meter balance is
\begin{align}
p_t^{\GI}
=p_t^{\EV}+p_t^{\ch,\BESS}-p_t^{\dch,\BESS}
+p_t^{\load}-p_t^{\PV}. \label{eq:meter}
\end{align}
The base case sets $p_t^{\load}=p_t^{\PV}=0$. In the sensitivity, the two terms are measured exogenous parameters and are not control variables. They affect TOU cost and the PD/NCD threshold trajectory only through $p_t^{\GI}$. They are excluded from DA and RT energy positions, all AS bids, the shared capability envelope, and all wholesale revenue. Retail settlement use the signed meter in \eqref{eq:meter}, including the adopted retail-price credit when $p_t^{\GI}<0$.

\subsection{DA/RT energy and AS settlement}
Let $\bar p_t^{\DA}$ and $\bar c_t^{x,\DA}$ be cleared DA energy and AS awards
already mapped to the 15-min model clock; their stored values are zero whenever
no award is active. The RT formulation fixes these arrays rather than
re-optimizing them:
\begin{align}
p_t^{\DA}&=\bar p_t^{\DA},&
p_t^{\actual}&=p_t^{\DA}+p_t^{\RT}, \label{eq:actual_energy}\\
c_t^{x,\DA}&=\bar c_t^{x,\DA},&
c_t^{x,\actual}&=c_t^{x,\DA}+c_t^{x,\RT}\ge0. \label{eq:actual_as}
\end{align}
Here $p_t^{\RT}$ and $c_t^{x,\RT}$ are signed settlement adjustments. A
negative AS adjustment records a reduction or repurchase of part of the DA
position, but the physically offered quantity $c_t^{x,\actual}$ cannot be
negative. This algebraic recourse is consistent with CAISO's separation of DA
and RT awards and its requirement that awarded AS remain deliverable; it does
not assert that CAISO publishes a negative physical capacity award
\cite{caisoMarketOperations2026} \cite{caisoNGRCommitment2020}. An explicit binary
award mask is unnecessary in the mathematics because it can be absorbed into
the cleared arrays $\bar p_t^{\DA}$ and $\bar c_t^{x,\DA}$.

\subsection{DA/RT award resolution and price normalization}
\label{sec:price_resolution}
The implementation keeps one 15-min computational clock while preserving the
different native market resolutions. CAISO DA energy and AS awards are hourly,
whereas the modeled RT binding awards are interval-specific
\cite{caisoMarketInstruments2026}. For the four intervals
$\calT_h$ belonging to hour $h$, the submitted DA variables therefore satisfy
\begin{align}
p_t^{\DA}&=p_h^{\DA},&
c_t^{x,\DA}&=c_h^{x,\DA},
\qquad t\in\calT_h,\quad |\calT_h|=4. \label{eq:da_hour_blocks}
\end{align}
They are represented four times only to align with the physical model; they
remain one hourly award. The signed $p_t^{\RT}$ and $c_t^{x,\RT}$ decisions
may change every 15 min. Once the DA award has cleared, the repeated values in
\eqref{eq:da_hour_blocks} enter every RT optimization call as constants rather than being
re-optimized.

Let $\widetilde\lambda_t^s$ denote a native energy LMP in US\$/MWh and
$\widetilde\rho_t^{x,s}$ a native ASMP reported in US\$/MW. A native RT
ASMP is not an interval-integrated payment: the 15-min capacity settlement
retains a factor of $0.25$, as illustrated by CAISO's regulation-down and
non-spinning-reserve settlement rules \cite{caisoRTRegDownSettlement} \cite{caisoRTNonSpinSettlement}.
All objective, table, and price-panel coefficients are normalized to
a common US\$/kWh-equivalent basis:
\begin{align}
\lambda_t^{\DA}&=\frac{\widetilde\lambda_h^{\DA}}{1000~\mathrm{kWh/MWh}},
&&t\in\calT_h,\\
\lambda_t^{\RT}&=\frac{\widetilde\lambda_t^{\RT}}{1000~\mathrm{kWh/MWh}},
&& \label{eq:lmp_normalization}\\
\rho_t^{x,\DA}&=\frac{\widetilde\rho_h^{x,\DA}}{(1000~\mathrm{kW/MW})(1~\mathrm{h})},
&&t\in\calT_h, \label{eq:da_asmp_normalization}\\
\rho_t^{x,\RT}&=\frac{\widetilde\rho_t^{x,\RT}}
{(1000~\mathrm{kW/MW})(1~\mathrm{h})}.
&& \label{eq:rt_asmp_normalization}
\end{align}
The same hourly-equivalent conversion applies in both markets; award-update
frequency does not remove the duration multiplier. With $\dt=0.25$ h,
one hourly DA capacity award receives
\begin{align}
\sum_{t\in\calT_h}\dt\rho_t^{x,\DA}c_h^{x,\DA}
&=4(0.25~\mathrm{h})
\frac{\widetilde\rho_h^{x,\DA}c_h^{x,\DA}}
{(1000~\mathrm{kW/MW})(1~\mathrm{h})}\nonumber\\
&=\frac{\widetilde\rho_h^{x,\DA}c_h^{x,\DA}}{1000~\mathrm{kW/MW}},
\label{eq:da_capacity_payment_check}
\end{align}
whereas a single RT binding award receives
\begin{align}
\dt\rho_t^{x,\RT}c_t^{x,\RT}
&=\dt\frac{\widetilde\rho_t^{x,\RT}}
{(1000~\mathrm{kW/MW})(1~\mathrm{h})}c_t^{x,\RT}\nonumber\\
&=0.25\frac{\widetilde\rho_t^{x,\RT}c_t^{x,\RT}}{1000~\mathrm{kW/MW}}.
\label{eq:rt_capacity_payment_check}
\end{align}
Thus the DA ASMP is \emph{not} divided by four; the hourly award is instead
held constant over its four physical intervals. The RT ASMP is likewise
not divided by the 15-min duration; the quarter-hour payment retains that
duration explicitly. Energy
LMPs require only the MWh-to-kWh conversion because their settlement already
contains $\dt$. The same indexing rule applies to energy: an hourly DA LMP and
its hourly $p_h^{\DA}$ award are repeated on the 15-min model clock, whereas
both the RT LMP and signed $p_t^{\RT}$ adjustment remain interval-specific.
Consequently the four stored DA rows are an accounting representation of one
hourly award, not four independently optimized DA awards.

With $P_{t,\max}^{\EV}=\sum_i\bar P_i b_{t,i}$, the common capability bounds are
\begin{align}
\bar P_t^{\mathrm{up}}&=P^{\BESS}+B_t^{\EV}, \label{eq:upbound}\\
\bar P_t^{\mathrm{down}}
&=\max\{P^{\BESS}-B_t^{\EV}+P_{t,\max}^{\EV},0\}. \label{eq:downbound}
\end{align}
The nonnegative clipping in \eqref{eq:downbound} is necessary when baseline EV
power exceeds the instantaneous BESS-plus-connected-EV absorption capability.
The EV--BESS energy position and actual AS share these bounds:
\begin{align}
\pos{p_t^{\actual}}+
c_t^{\RU,\actual}+c_t^{\SP,\actual}+c_t^{\NSP,\actual}
&\le\bar P_t^{\mathrm{up}}, \label{eq:up_envelope}\\
\pos{-p_t^{\actual}}+c_t^{\RD,\actual}
&\le\bar P_t^{\mathrm{down}}. \label{eq:down_envelope}
\end{align}
The DA position and signed RT adjustment are likewise bounded by the
corresponding interval capability limits. These limits are identical with and
without passive PV/building because those profiles are not dispatchable.

Expected AS deployment gives the EV--BESS market-facing obligation
\begin{align}
q_t^{\WM}={}&p_t^{\actual}
+\alpha_t^{\RU}c_t^{\RU,\actual}
-\alpha_t^{\RD}c_t^{\RD,\actual}\nonumber\\
&+\alpha_t^{\SP}c_t^{\SP,\actual}
+\alpha_t^{\NSP}c_t^{\NSP,\actual}. \label{eq:qwm}
\end{align}
The net-flow direction constraints below link the market obligation to the device-level accounting channels.

In \eqref{eq:qwm}, $q_t^{\WM}$ is the actual EV--BESS energy obligation,
$\alpha_t^x$ is the expected activated fraction of AS product $x$, and the
minus sign on RD represents additional absorption. This same obligation is
linked to the EV and BESS wholesale channels in \ref{app:complete_model}.

The reference case uses minimum duration requirement $T^{\mathrm{DU}}=0.5$ h for all AS products as a modeling
assumption, not as a claim that every CAISO product has identical deployment
and certification rules. The attenuation-factor sensitivity uses CAISO's hourly
regulation-up and regulation-down attenuation factors
\cite{caisoAttenuation2025}; spinning and non-spinning reserve retain explicit
assumptions unless a separate event-level deployment data set is supplied.

\subsection{Retail costs}
The TOU term is defined as:
\begin{align}
C_k^{\TOU}
=\sum_{t\in\calH_k}C_t^{\TOU}p_t^{\GI}\dt. \label{eq:tou}
\end{align}
Thus $p_t^{\GI}<0$ receives a retail-price export credit under the adopted
tariff convention.

For demand category $q\in\{\PD,\NCD\}$, let $\calT_q$ be its tariff-eligible
intervals and let $M_k^{q,\hist}$ be the largest eligible import executed
before call $k$. The avoidable increase in the monthly peak is represented by
the epigraph variable $z_k^q$:
\begin{align}
z_k^q&\ge0, \qquad
z_k^q\ge p_t^{\GI}-M_k^{q,\hist},
\quad t\in\calH_k\cap\calT_q, \label{eq:demand_epi}\\
C_k^{q,\mathrm{inc}}&=C^q z_k^q. \label{eq:demand_inc}
\end{align}
After executing interval $k$, the exogenous indicator $m_q(k)$ equals one when
the interval belongs to category $q$ and zero otherwise, so the carried peak
updates as
\begin{align}
M_{k+1}^{q,\hist}
=\max\{M_k^{q,\hist},\,m_q(k)p_k^{\GI,\mathrm{executed}}\}. \label{eq:threshold_update}
\end{align}
Thus the objective charges only a newly established billing peak; previously
incurred demand cost is a constant and cannot influence the current decision.
The equivalent absolute and incremental forms are derived in \ref{app:closed_loop}.

\subsection{Wholesale revenue and objective}
Wholesale energy revenue, AS capacity revenue, and EV charging-service revenue
are, respectively,
\begin{align}
R_k^{\WM,\mathrm{energy}}
={}&\sum_{t\in\calH_k}\dt
\left(\lambda_t^{\DA}p_t^{\DA}
+\lambda_t^{\RT}p_t^{\RT}\right)+R_k^{\mathrm{deploy}}, \label{eq:wm_energy}\\
R_k^{\WM,\mathrm{capacity}}
={}&\sum_{t\in\calH_k}\dt\sum_{x\in\calK}
\left(\rho_t^{x,\DA}c_t^{x,\DA}
+\rho_t^{x,\RT}c_t^{x,\RT}\right), \label{eq:wm_capacity}\\
R_k^{\EV}
={}&\sum_{t\in\calH_k}C^{\EV}p_t^{\EV}\dt. \label{eq:ev_revenue}
\end{align}
$R_k^{\EV}$ is the payment from EV drivers for delivered charging energy at rate $C^{\EV}$. In \eqref{eq:wm_energy}, $\lambda_t^{\DA}$ and $\lambda_t^{\RT}$ settle the fixed
DA energy position and signed RT adjustment. In \eqref{eq:wm_capacity}, $\rho_t^{x,\DA}$ and $\rho_t^{x,\RT}$ settle the corresponding AS capacities. $R_k^{\mathrm{deploy}}$ prices the expected AS deployment energy in
\eqref{eq:qwm} at the RT LMP:
\begin{align}
R_k^{\mathrm{deploy}}
={}&\sum_{t\in\calH_k}\dt\lambda_t^{\RT}
\bigg[\alpha_t^{\RU}c_t^{\RU,\actual}
-\alpha_t^{\RD}c_t^{\RD,\actual}\nonumber\\
&+\alpha_t^{\SP}c_t^{\SP,\actual}
+\alpha_t^{\NSP}c_t^{\NSP,\actual}\bigg].
\label{eq:wm_deployment_revenue}
\end{align}
The sign difference is essential: RD deployment imports energy and therefore
subtracts RT energy value under the adopted positive-delivery convention,
whereas RU, SP, and NSP add upward delivery.

At MPC call $k$, the complete active-horizon problem is
\begin{align}
\min_{\boldsymbol{x}_k\in\mathcal{X}(\mathcal{I}_k)}
J_k={}&
\sum_{t\in\calH_k}C_t^{\TOU}p_t^{\GI}\dt
+\sum_{q\in\{\PD,\NCD\}}C^qz_k^q \nonumber\\
&-R_k^{\EV}
-R_k^{\WM,\mathrm{energy}}
-R_k^{\WM,\mathrm{capacity}}, \label{eq:objective}
\end{align}
subject to the key coupling constraints above and the complete feasible set in \ref{app:complete_model}. The first line is retail energy plus the incremental monthly demand charges; the second line subtracts EV charging-service revenue and the two wholesale revenue accounts. There is no DA-tracking penalty. The economic consequence of deviating from DA is already represented by the DA/RT two-settlement terms.

\section{Forecast updaters, horizon, and oracle logic}
\label{sec:oracle_logic}

\subsection{Perfect-information and Persistence updaters}
Once a vehicle arrives, its realized session information is used. For sessions
that have not arrived, Perfect information supplies the realized future session set within
the active horizon, while Persistence constructs the future set from recent
charging behavior. Electricity prices are exogenous processed inputs in both
cases; ``Perfect'' refers only to EV information within the rolling 24-h
window, not to price forecasting or full-period knowledge.

\subsection{Why 24-h Perfect information need not maximize billing-period value}
Let $\omega$ denote the complete realized billing-period data and
$\mathcal{X}(\omega)$ the common full-period feasible set. A true oracle computes
\begin{align}
J_{\mathrm{oracle}}^\star(\omega)
=\min_{\boldsymbol{x}\in\mathcal{X}(\omega)}
J_{\calT}(\boldsymbol{x};\omega). \label{eq:oracle}
\end{align}
For every feasible policy $\pi$,
\begin{align}
J_{\mathrm{oracle}}^\star\le J^\pi
\quad\Longleftrightarrow\quad
V_{\mathrm{oracle}}^\star\ge V^\pi,\qquad V=-J. \label{eq:oracle_dominance}
\end{align}
The inequality direction differs for cost and revenue. These equations define a conditional theoretical reference. The monthly site comparison uses the two 24-h forecast updaters; the separate weekly benchmark reports its DA-envelope difference explicitly and is not asserted to bound every forecast-driven schedule.

A 24-h controller using updater $\pi$ instead selects the decision trajectory
$\boldsymbol{x}_{k,H}^{\pi}$ over intervals $k,\ldots,k+H-1$ by minimizing the
explicit truncated economic objective
\begin{align}
\boldsymbol{x}_{k,H}^{\pi}
&\in\arg\min_{\boldsymbol{x}\in\mathcal X_k^\pi}
\Bigg\{\sum_{\tau=k}^{k+H-1}C_\tau^{\TOU}p_\tau^{\GI}\dt
+\sum_{q\in\{\PD,\NCD\}}C^qz_k^q\nonumber\\
&-R_{k,H}^{\EV}-R_{k,H}^{\WM,\mathrm{energy}}
-R_{k,H}^{\WM,\mathrm{capacity}}\Bigg\}, \label{eq:mpc}
\end{align}
where $H=96$ intervals and the three revenue terms are the restrictions of
\eqref{eq:wm_energy}--\eqref{eq:ev_revenue} to that horizon. It executes only
$\boldsymbol{x}_k^\pi$ and repeats. No exact cost-to-go from $k+H$ to the end
of the billing period appears in \eqref{eq:mpc}. Consequently, more
accurate information inside the truncated horizon does not imply globally
better closed-loop value \cite{rawlings2017mpc}.

The mechanism resembles a greedy algorithm. Each optimization call selects a locally
attractive first action from incomplete future information. That action changes
the next state and therefore the feasible set of every later call. Four state
channels matter here:
\begin{enumerate}
  \item BESS SOC and daily throughput usage;
  \item monthly PD and NCD thresholds;
  \item EV energy remaining before departure; and
  \item submitted DA energy and AS commitments that cannot be undone.
\end{enumerate}
A Persistence forecast can accidentally preserve a more valuable out-of-horizon state than a 24-h Perfect-information optimization. A 1\% mixed-integer linear programming (MILP) optimality gap can also select different near-optimal first actions and amplify their differences through the cascade, but this is a secondary numerical mechanism. The valid requirement is that the full-period oracle dominate every feasible executed trajectory; 24-h Perfect information versus Persistence remains an empirical
controller comparison.


\section{Case study and implementation}
\label{sec:case}
\label{sec:case_study}

\subsection{Modeling assumptions}
The study adopts the following assumptions so that the reported economic
comparison is not mistaken for a complete CAISO participation study.
\begin{enumerate}
  \item DA/RT energy prices, AS prices, and retail prices are known exogenous
  inputs over each active optimization horizon; price-forecast error is outside
  the present forecast comparison.
  \item When an EV arrives, its departure, requested energy, and charger limit
  become known to the controller. The Perfect-information updater also
  knows realized but not-yet-arrived sessions within the 24-h horizon, whereas
  Persistence does not.
  \item Submitted energy and AS quantities are assumed to clear. DA awards are
  fixed in RT, and signed RT recourse represents an incremental purchase or
  sale around the DA position. Qualification, minimum-bid, telemetry, and
  performance-score requirements are not modeled \cite{caisoAncillary2026} \cite{caisoDER2026}.
  \item No AGC signal is simulated. RU/RD use the hourly CAISO attenuation
  factors for the corresponding quarter, repeated over the four 15-min intervals;
  SP/NSP retain the assumed factor 0.2. These are expected-energy approximations,
  not measured realization-specific deployment fractions \cite{caisoAttenuation2025}.
  \item PV generation and building demand are measured passive PCC quantities.
  They affect retail cost but provide neither wholesale energy nor AS capacity.
  \item Retail export is credited at the same modeled TOU coefficient as import.
  This symmetric convention is an analytical assumption, not an SDG\&E export
  tariff \cite{sdgeLargeCommercial2025}.
  \item The wholesale baseline-relative position and the signed retail meter
  are coupled analytical accounts. Their simultaneous remuneration is a study
  assumption, not a determination that a particular interconnection or tariff
  permits both credits. Implementation would require verification of resource
  registration, metering boundaries, eligible settlements, and any rules
  preventing duplicate compensation \cite{caisoDER2026}.
  \item Each reported site is represented by one aggregate PCC. Feeder congestion,
  individual branch-circuit limits, V2G export, and detailed BESS degradation are
  outside the present scope.
\end{enumerate}

\subsection{Site-specific UCSD workplace-charging portfolios}
The reported case study selects two mapped configurations from the broader processed UCSD archive. NTPLL combines seven building-meter channels across six mapped structures, three PV meters, and 34 EVSEs at Scholars Parking. Center Hall combines one building meter, the nearby Osler/South Parking PV meter, and 38 EVSEs at South Parking. Both cases use the same BESS rating and optimization equations, so differences arise from measured site demand, PV, and Level 2 EV charging sessions rather than a retuned storage asset. The dominant charging pattern is daytime workplace charging: many users arrive in the morning, remain connected for several hours, and depart during the afternoon or evening.

\subsection{Physical and economic parameters}
Table~\ref{tab:case_parameters} summarizes the principal implemented settings. The BESS uses the geometric efficiency split in \eqref{eq:soc}; hence $\sqrt{\gamma}$ applies to charging and $1/\sqrt{\gamma}$ to discharging. The 0.5-h AS-duration guard is conservative and common to the upward and downward reserve equations. The final billing-period SOC is fixed at 50\% so no policy can increase reported value by borrowing energy from beyond the study boundary.

\begin{table*}[!t]
\centering
\caption{Principal case-study parameters used by the validated implementation.}
\label{tab:case_parameters}
\small
\begin{tabularx}{\textwidth}{@{}p{0.29\textwidth}p{0.22\textwidth}X@{}}
\toprule
Parameter & Value & Interpretation \\
\midrule
Time step and controller horizon & 0.25 h; 24 h & One decision is executed before the horizon advances. \\
BESS power and energy & 250 kW; 332 kWh & Common stationary-storage limits. \\
BESS SOC and efficiency & 0.05--0.95; $\gamma=0.90$ & Energy state and round-trip factor. \\
Initial/final billing-period SOC & 0.50 / 0.50 & Common monthly boundary conditions; historical EV dispatch is carried between months. \\
Nominal EV charger power & 6.656 kW/port & Further limited by the measured interval envelope. \\
EV user payment & \$0.40/kWh & Identical aggregate service gives identical monthly EV revenue. \\
Summer TOU (on/off/super-off) & 0.82500 / 0.21421 / 0.12071 \$/kWh & June--October; frozen study coefficients. \\
Winter TOU (on/off/super-off) & 0.55300 / 0.13478 / 0.11278 \$/kWh & November--May; frozen study coefficients. \\
PD: summer / winter & 3.05 / 0.63 \$/kW & New monthly on-peak maximum. \\
NCD: both seasons & 15.38 \$/kW & New monthly maximum over all intervals. \\
RU attenuation range: Q1/Q2/Q3 & 0.024--0.177 / 0.041--0.204 / 0.048--0.216 & Hour-specific CAISO values; Q1 applies to March, Q2 to April--June, Q3 to July. \\
RD attenuation range: Q1/Q2/Q3 & 0.299--0.518 / 0.243--0.605 / 0.227--0.558 & Hour-specific CAISO values, not one constant per quarter. \\
SP/NSP attenuation factors & 0.2 / 0.2 & Explicit assumptions; not inferred from RU/RD data. \\
AS duration requirement & 0.5 h & SOC guard for reserved response. \\
MILP tolerance & 1\% relative gap & Accepted for every detailed 24-h optimization. \\
\bottomrule
\end{tabularx}
\end{table*}

Retail TOU prices are represented by the composite coefficients used in the
case study. The modeled large-commercial schedule distinguishes summer
(June--October) and winter (November--May) energy periods; PD is restricted to
the tariff-defined on-peak window, whereas NCD uses the maximum eligible
billing-cycle demand \cite{sdgeLargeCommercial2025} \cite{sdgeDemandGuide2026}. The
study treats export symmetrically in the signed TOU expression. The numerical
coefficients are the frozen SDG\&E-based study tariff, not a reconstruction of
every effective-date tariff revision in 2025. In the implemented calendar,
16:00--21:00 is on-peak every day. On weekdays, 00:00--06:00 and 10:00--14:00
are super-off-peak; on weekends and modeled US holidays, the off-peak windows
are 14:00--16:00 and 21:00--24:00, and the remaining non-on-peak hours are
super-off-peak. This common schedule isolates measured-site and forecast effects.
CAISO's hourly RU/RD multipliers describe anticipated SOC effects and are
updated quarterly from historical use; applying them as expected deployed-energy
factors is a modeling approximation \cite{caisoAttenuation2025}. 

\subsection{Scenario design}
The site comparison contains four resource/information stages:
\begin{description}
  \item[Building only:] reconstructed gross building demand is billed without PV, EV charging, BESS operation, or wholesale participation;
  \item[Building + PV:] measured PV is subtracted once, recovering the native net-meter trajectory; this is a passive retail calculation rather than an optimization;
  \item[24-h Persistence:] historical session information forms the unarrived-EV forecast, and the RT problem replaces it with realized arrivals as they become observable;
  \item[24-h Perfect information:] realized EV sessions inside the 24-h window are supplied in advance, while the control horizon and all other assumptions remain unchanged.
\end{description}
The last two stages both include building, PV, EV, and BESS. They are the like-for-like controller comparison and differ only in EV information. The first two stages contain no controllable resource and are evaluated directly from the measured meter data and tariff; their differences must not be interpreted as controller revenue.

\section{Financial accounting and validation}
\label{sec:accounting}

\subsection{Executed financial statement}
Every MPC optimization call produces a plan, but financial evaluation uses only the implemented first interval. For an evaluation set $\mathcal T_E$, the reconciled net value is
\begin{align}
V(\mathcal T_E)={}&R^{\EV}+R^{\WM,\mathrm{energy}}
+R^{\WM,\mathrm{capacity}}\nonumber\\
&+C^{\TOU}+C^{\PD}+C^{\NCD},
\label{eq:financial_identity}
\end{align}
where costs are stored as negative values. EV revenue is a service-payment account and is identical when the total delivered energy is identical. WM revenue is reconstructed from the executed DA position, signed RT adjustment, AS capacity, and expected deployment. TOU cost is accumulated interval by interval from the signed meter. PD and NCD are charged only when the executed meter establishes a new billing-period maximum.

Daily demand-charge attribution is an accounting decomposition of a monthly charge. A large negative value on one day indicates that the day established a new threshold; it is not a separate daily tariff. Monthly and cumulative figures therefore carry the economic conclusion, while daily figures explain the mechanism. The implementation accepts a maximum \$0.02 discrepancy between the daily rounded components and their independently computed total.

\subsection{Validation protocol}
The validation first confirms complete, unique time axes and matched scenario
inputs. It then tests optimizer completion, physical equations, financial
identities, EV service, and boundary states. Finally, interval trajectories are retained for every operating day, and representative days are rendered so numerical feasibility can be compared with the prices and bids that produced it.

\begin{table*}[!t]
\centering
\caption{Formal validation hierarchy and acceptance criteria.}
\label{tab:validation_protocol}
\scriptsize
\begin{tabularx}{\textwidth}{@{}p{0.27\textwidth}p{0.30\textwidth}X@{}}
\toprule
Layer & Criterion & Purpose \\
\midrule
Solver completeness & one DA and 96 RT records per operating day and forecast; accepted status; gap $\le1\%$; no TimeLimit & Excludes silent truncation and missing intervals.\\
Physical identities & Meter, settlement, and actual up/down residuals $\le10^{-6}$ kW; actual AS $\ge-10^{-5}$ kW & Confirms the implemented dispatch matches the model.\\
EV service & Billing-period EV-energy spread across matched forecasts $\le10^{-4}$ kWh & Prevents financial gains from unequal charging service.\\
Financial accounting & Daily component sum agrees with total within \$0.02 & Allows independent cent rounding but no missing account.\\
Terminal state & Final billing-period SOC equals 0.5 within $10^{-6}$ & Prevents end-of-study energy borrowing.\\
Scenario comparability & common site data, prices, BESS, tariff, and terminal conditions & Isolates the EV-information update within each site.\\
Graphical audit & Saved interval trajectories and representative daily figures & Links numerical checks to physical trajectories.\\
\bottomrule
\end{tabularx}
\end{table*}

\subsection{Computational acceptance}
The March--July set contains 20 site/month/forecast cases and 58,752 recorded
RT steps (153 days, two sites, and two forecast updaters). The saved validation
summaries report accepted optimal status at the requested 1\% relative-gap
setting and no time-limit event. This tolerance is not an exact global-optimality
certificate. Matched daily EV service differs by less than $10^{-4}$ kWh, and
every monthly endpoint satisfies SOC $=0.5$. All simulations use Python and
CVXPY with Gurobi as the optimization solver.

\section{Results}
\label{sec:results}

\subsection{Site power scale and measurement boundary}
\label{sec:site_scale_results}

Figure~\ref{fig:site_maps} identifies the two UCSD configurations used in the
formal experiment. NTPLL represents a mixed-use neighborhood with six mapped
structures, seven building-meter channels, three distributed PV meters, and 34
EVSEs at Scholars Parking. Center Hall represents a smaller academic building
paired with the nearby Osler/South Parking PV and 38 EVSEs. Feature coordinates
were obtained from OpenStreetMap and cross-checked against the official UCSD
NTPLLN neighborhood and main-campus maps \cite{openstreetmap2026}
\cite{ucsdNTPLLNMap2021} \cite{ucsdCampusMap2025}.

\begin{figure*}[!t]
  \centering
  \includegraphics[width=0.94\textwidth]{fig_03_ucsd_site_location_maps.png}
  \caption{Mapped NTPLL and Center Hall configurations. Symbols identify the
  building-meter, PV, and EVSE groups used by the site preprocessing. The map
  establishes data membership; it does not imply a feeder power-flow model.}
  \label{fig:site_maps}
\end{figure*}

Table~\ref{tab:site_power_specs} reports the corresponding power scales. The PV
capacity proxy is a data-derived observed-capacity field rather than a certified
nameplate rating. The aggregate charger limit is the sum of modeled per-port
limits, whereas the illustrative June EV peak is the maximum coincident measured charging
power. These distinctions prevent installed capability from being confused with
realized demand.

\begin{table*}[!t]
\centering
\caption{Measured and modeled site power specifications. Peaks are calculated
from the June 2025 15-min data; Fig.~\ref{fig:site_power_specs} additionally shows all five months. PV proxy values are not certified nameplate
ratings.}
\label{tab:site_power_specs}
\scriptsize
\setlength{\tabcolsep}{3.5pt}
\begin{tabular}{@{}lrrrrrrrr@{}}
\toprule
Site & \shortstack{Building\\channels} & \shortstack{PV\\meters} & EVSE &
\shortstack{Charger\\limit (kW)} & \shortstack{EV peak\\(kW)} &
\shortstack{PV proxy / peak\\(kW)} & \shortstack{Gross building /\\native net peak (kW)} &
\shortstack{BESS\\(kW/kWh)}\\
\midrule
NTPLL & 7 & 3 & 34 & 226.304 & 92.993 & 166.930 / 163.208 & 1,398.239 / 1,263.811 & 250 / 332\\
Center Hall & 1 & 1 & 38 & 252.928 & 101.107 & 206.150 / 207.714 & 295.531 / 102.994 & 250 / 332\\
\bottomrule
\end{tabular}
\end{table*}

The sites are intentionally not normalized to the same passive-load scale.
NTPLL's gross building peak is approximately 4.7 times that of Center Hall,
whereas the same 250-kW/332-kWh BESS is used in both cases. Center Hall also has
a much larger PV-to-building ratio. Figure~\ref{fig:site_power_specs} makes
these relative scales visible and explains why equal absolute storage and
charger ratings need not produce equal tariff or forecast effects.

\begin{figure*}[!t]
  \centering
  \includegraphics[width=0.96\textwidth,height=0.83\textheight,keepaspectratio]{fig_04_site_system_power_specifications_mar_jul_2025.png}
  \caption{March--July 2025 site power specifications, with one row per month. Installed/model limits and observed
  coincident peaks are shown separately; the common BESS therefore occupies a
  much larger fraction of the Center Hall site envelope than of NTPLL's.}
  \label{fig:site_power_specs}
\end{figure*}

The native building channels are net meters that already contain embedded PV.
The preprocessing identity is
\begin{equation}
\begin{aligned}
p_t^{\load,\mathrm{gross}}
&=p_t^{\mathrm{meter,native}}+p_t^{\PV},\\
p_t^{\load,\mathrm{gross}}-p_t^{\PV}
&=p_t^{\mathrm{meter,native}}.
\end{aligned}
\label{eq:site_net_meter_identity}
\end{equation}
Thus the building-only counterfactual removes PV by reconstructing gross load,
whereas the building-plus-PV case returns the original net-meter trajectory.
PV is subtracted exactly once in the optimized cases.

\subsection{Monthly value across resource stages}
\label{sec:site_monthly_value}

Table~\ref{tab:site_monthly_results} and
Fig.~\ref{fig:site_resource_financial} compare three resource stages. Building
only and building plus PV are deterministic retail-billing counterfactuals; no
optimizer or wholesale account is used. The full stage adds EV charging, the
BESS, and the wholesale model and is solved separately with the Persistence and
Perfect EV forecast updaters. Costs are shown as negative values, so a less
negative net value is economically better.

\begin{table*}[!t]
\centering
\caption{March--July 2025 monthly net value (US\$). Passive stages contain no EV-service or wholesale revenue. Both full-portfolio columns include building, PV, EV and BESS. Positive Perfect minus Persistence favors Perfect. Monthly values are summed before final rounding.}
\label{tab:site_monthly_results}
\scriptsize
\setlength{\tabcolsep}{5pt}
\begin{tabular}{@{}llrrrrr@{}}
\toprule
Month & Site & Building only & Building + PV & Persistence & Perfect & Perfect $-$ Persistence\\
\midrule
Mar & Center Hall & -20,443.60 & -12,177.52 & -2,662.21 & -2,404.86 & 257.35 \\
Mar & NTPLL & -152,091.48 & -147,256.36 & -135,867.36 & -135,662.39 & 204.96 \\
Apr & Center Hall & -22,999.15 & -12,795.80 & -4,321.27 & -4,856.52 & -535.26 \\
Apr & NTPLL & -172,949.64 & -165,048.57 & -154,729.66 & -154,532.50 & 197.16 \\
May & Center Hall & -22,438.82 & -12,818.38 & -4,404.27 & -4,657.27 & -253.00 \\
May & NTPLL & -195,582.36 & -188,191.12 & -176,413.15 & -176,467.73 & -54.58 \\
Jun & Center Hall & -27,603.53 & -15,752.16 & -6,880.74 & -7,106.78 & -226.04 \\
Jun & NTPLL & -224,173.75 & -215,474.03 & -205,691.93 & -205,657.63 & 34.30 \\
Jul & Center Hall & -28,823.60 & -15,770.65 & -7,896.91 & -7,437.80 & 459.11 \\
Jul & NTPLL & -210,201.20 & -201,566.24 & -187,960.86 & -188,110.08 & -149.21 \\
\bottomrule
\end{tabular}
\end{table*}

Table~\ref{tab:site_monthly_decomposition} supplements the net-value comparison
with the complete WM, EV-service, TOU, PD, and NCD accounts. Both forecast
rows represent the full portfolio; the passive rows contain no EV or WM receipts.
\begin{table*}[!t]
\centering
\caption{Monthly financial decomposition, March--July 2025 (US\$). Costs are negative. Persistence and Perfect denote the full building--PV--EV--BESS portfolio. Values are accumulated before rounding; small differences between displayed components and totals are rounding only.}
\label{tab:site_monthly_decomposition}
\footnotesize
\setlength{\tabcolsep}{3pt}
\begin{tabular}{@{}lllrrrrrr@{}}
\toprule
Month & Site & Case & WM revenue & EV revenue & TOU & PD & NCD & Net value\\
\midrule
Mar & NTPLL & Building only & 0.00 & 0.00 & -132,185.26 & -744.53 & -19,161.69 & -152,091.48 \\
 &  & Building + PV & 0.00 & 0.00 & -127,897.71 & -733.93 & -18,624.72 & -147,256.36 \\
 &  & Persistence & 4,816.76 & 2,743.65 & -126,110.87 & -681.43 & -16,635.47 & -135,867.36 \\
 &  & Perfect & 4,802.83 & 2,743.65 & -126,028.01 & -676.07 & -16,504.79 & -135,662.39 \\
\addlinespace
Mar & Center Hall & Building only & 0.00 & 0.00 & -15,769.41 & -135.76 & -4,538.42 & -20,443.60 \\
 &  & Building + PV & 0.00 & 0.00 & -10,511.40 & -63.31 & -1,602.80 & -12,177.52 \\
 &  & Persistence & 5,116.40 & 4,095.10 & -8,829.50 & -58.21 & -2,986.00 & -2,662.21 \\
 &  & Perfect & 4,918.19 & 4,095.10 & -9,140.39 & -72.15 & -2,205.62 & -2,404.86 \\
\addlinespace
Apr & NTPLL & Building only & 0.00 & 0.00 & -149,191.49 & -826.94 & -22,931.21 & -172,949.64 \\
 &  & Building + PV & 0.00 & 0.00 & -143,660.75 & -826.81 & -20,561.01 & -165,048.57 \\
 &  & Persistence & 3,779.04 & 3,728.10 & -142,745.83 & -766.19 & -18,724.78 & -154,729.66 \\
 &  & Perfect & 3,657.69 & 3,728.10 & -142,633.74 & -756.97 & -18,527.57 & -154,532.50 \\
\addlinespace
Apr & Center Hall & Building only & 0.00 & 0.00 & -18,010.02 & -148.11 & -4,841.02 & -22,999.15 \\
 &  & Building + PV & 0.00 & 0.00 & -11,118.10 & -64.28 & -1,613.42 & -12,795.80 \\
 &  & Persistence & 3,955.35 & 4,033.43 & -9,656.04 & -69.38 & -2,584.64 & -4,321.27 \\
 &  & Perfect & 3,720.54 & 4,033.43 & -10,421.91 & -73.00 & -2,115.58 & -4,856.52 \\
\addlinespace
May & NTPLL & Building only & 0.00 & 0.00 & -170,324.01 & -890.78 & -24,367.57 & -195,582.36 \\
 &  & Building + PV & 0.00 & 0.00 & -165,249.29 & -874.57 & -22,067.25 & -188,191.12 \\
 &  & Persistence & 3,934.36 & 3,429.84 & -163,980.91 & -763.10 & -19,033.33 & -176,413.15 \\
 &  & Perfect & 3,942.71 & 3,429.84 & -164,147.86 & -771.48 & -18,920.93 & -176,467.73 \\
\addlinespace
May & Center Hall & Building only & 0.00 & 0.00 & -17,567.82 & -146.39 & -4,724.62 & -22,438.82 \\
 &  & Building + PV & 0.00 & 0.00 & -11,158.28 & -62.78 & -1,597.32 & -12,818.38 \\
 &  & Persistence & 4,102.79 & 3,887.64 & -9,826.42 & -64.02 & -2,504.26 & -4,404.27 \\
 &  & Perfect & 3,965.95 & 3,887.64 & -10,464.24 & -69.38 & -1,977.25 & -4,657.27 \\
\addlinespace
Jun & NTPLL & Building only & 0.00 & 0.00 & -198,917.47 & -3,751.36 & -21,504.92 & -224,173.75 \\
 &  & Building + PV & 0.00 & 0.00 & -192,285.25 & -3,751.36 & -19,437.42 & -215,474.03 \\
 &  & Persistence & 3,196.91 & 1,906.31 & -188,756.28 & -3,632.05 & -18,406.82 & -205,691.93 \\
 &  & Perfect & 3,204.24 & 1,906.31 & -188,773.88 & -3,611.56 & -18,382.74 & -205,657.63 \\
\addlinespace
Jun & Center Hall & Building only & 0.00 & 0.00 & -22,379.67 & -678.59 & -4,545.27 & -27,603.53 \\
 &  & Building + PV & 0.00 & 0.00 & -13,870.74 & -297.37 & -1,584.05 & -15,752.16 \\
 &  & Persistence & 3,407.56 & 2,644.74 & -10,415.85 & -344.47 & -2,172.72 & -6,880.74 \\
 &  & Perfect & 3,292.04 & 2,644.74 & -11,055.96 & -305.07 & -1,682.53 & -7,106.78 \\
\addlinespace
Jul & NTPLL & Building only & 0.00 & 0.00 & -187,877.10 & -3,694.44 & -18,629.66 & -210,201.20 \\
 &  & Building + PV & 0.00 & 0.00 & -180,272.44 & -3,523.93 & -17,769.86 & -201,566.24 \\
 &  & Persistence & 3,182.54 & 1,204.51 & -175,654.34 & -2,761.43 & -13,932.14 & -187,960.86 \\
 &  & Perfect & 3,186.08 & 1,204.51 & -175,783.56 & -2,761.43 & -13,955.67 & -188,110.08 \\
\addlinespace
Jul & Center Hall & Building only & 0.00 & 0.00 & -23,660.28 & -675.42 & -4,487.90 & -28,823.60 \\
 &  & Building + PV & 0.00 & 0.00 & -14,039.61 & -277.00 & -1,454.03 & -15,770.65 \\
 &  & Persistence & 3,311.41 & 2,758.68 & -11,849.06 & -330.34 & -1,787.60 & -7,896.91 \\
 &  & Perfect & 3,318.60 & 2,758.68 & -11,585.02 & -277.00 & -1,653.07 & -7,437.80 \\
\addlinespace
\bottomrule
\end{tabular}
\end{table*}

\begin{figure*}[!t]
  \centering
  \includegraphics[width=0.96\textwidth,height=0.83\textheight,keepaspectratio]{fig_05_site_resource_addition_financial_comparison_mar_jul_2025.png}
  \caption{March--July net-value and component changes (one row per month) as PV and then coordinated
  EV--BESS operation are added. Building-only and building-plus-PV bars are
  direct tariff reconstructions; only the full-portfolio bars are optimized.}
  \label{fig:site_resource_financial}
\end{figure*}

The value stack varies substantially over the five months. PV improves the
passive net value by \$4,835--\$8,700 per month at NTPLL and
\$8,266--\$13,053 at Center Hall. Relative to building plus PV, the full
Persistence portfolio adds \$9,782--\$13,605 at NTPLL and
\$7,874--\$9,515 at Center Hall. These increments include EV charging-service
receipts and WM revenue as well as changes in retail cost; they are not pure
controller savings or investment returns. Nevertheless, the Persistence full
portfolio improves on its passive predecessor in every site/month pair.

Figure~\ref{fig:site_resource_financial} separates two mechanisms that a single
summer-month total would conceal. At NTPLL, Persistence WM revenue falls from
\$4,817 in March to \$3,183 in July, while the negative TOU account grows from
\$126,111 in March to \$188,756 in June before declining to \$175,654 in July.
The summer tariff change therefore coincides with a different passive-load and
EV-service trajectory; it cannot be interpreted as an isolated tariff experiment.
At Center Hall, the smaller passive bill makes the same order of wholesale
receipts much more visible in net value. Its Persistence value moves from
$-\$2,662$ in March to $-\$7,897$ in July even though each month still improves on
building plus PV. Absolute profit, incremental portfolio value, and forecast
value answer different questions and must not be conflated.

\subsection{Forecast-updater comparison}
\label{sec:site_forecast_results}

Forecast value changes with both site and month (Table~\ref{tab:site_monthly_results}).
Perfect improves NTPLL value in March, April, and June by \$34--\$205, but is
\$55 worse in May and \$149 worse in July. At Center Hall it improves March
and July by \$257 and \$459, respectively, but is \$226--\$535 worse in April--June.
Thus neither a single site's ranking nor a single month's ranking generalizes
to all five billing periods.

June illustrates the competing accounts. At NTPLL, the \$34.30 improvement
comprises \$7.33 additional WM revenue and reductions of \$20.49 in PD and
\$24.08 in NCD, offset by \$17.60 additional TOU cost. At Center Hall, Perfect
reduces PD and NCD by \$39.40 and \$490.19, but loses \$115.52 of WM revenue
and incurs \$640.11 more TOU cost, leaving net value \$226.04 lower. In July,
NTPLL's \$3.54 WM improvement is outweighed by \$129.22 additional TOU cost
and \$23.53 additional NCD cost. EV-service receipts are identical
within each matched site/month pair, so unequal charging service does not
explain these differences.

This result is consistent with the rolling formulation in
Section~\ref{sec:oracle_logic}. Perfect information is exact only inside each
24-h window; it is not a billing-period oracle. Executed actions carry SOC,
remaining EV energy, demand thresholds, and irreversible DA commitments into
future optimizations. A locally superior trajectory can therefore create a
less valuable state beyond the current horizon. The site comparison should be
read as the value of the forecast updater inside the implemented controller,
not as a theorem that Perfect information must dominate Persistence.



\subsection{Daily market-settlement audit}
\label{sec:site_daily_audit}

Figures~\ref{fig:ntpll_daily_audit}--\ref{fig:center_daily_audit_later}
show Perfect on July 1 and 15 at both sites. Each site's corresponding panels
use identical limits across dates. Prices use \$/kWh, with a separate TOU
axis. In panel (a), the total expected marginal product price for AS product $x$ and
market $s\in\{\DA,\RT\}$ is
\begin{equation}
\pi_t^{x,s}=\rho_t^{x,s}+\sigma_x\alpha_t^x\lambda_t^{\RT},
\qquad \sigma_{\RD}=-1,\quad \sigma_{\RU,\SP,\NSP}=1.
\label{eq:expected_product_price}
\end{equation}
Here $\rho$ is the normalized capacity coefficient and the second term is the
expected signed deployment-energy contribution per kW-hour of capacity. Both
DA and RT AS deployment use the RT LMP. Energy curves instead show the native
DA/RT LMP coefficients. This is a marginal price diagnostic, not an additional
settlement account or a guaranteed profit after retail cost and physical
opportunity costs. Panel (b) removes the deployment component to show capacity
prices separately. Panels (c) and (d) retain separate DA/RT bars; the black
line in (d) is precisely $q_t^{\WM}$ in \eqref{eq:qwm}. Signed buy-backs can
make separately stacked accounting bars exceed the physical lines in (c);
deliverability applies to the actual combined positions, not their gross stack.

At NTPLL on July 1, the BESS takes in 157.48 kWh before 06:00 and reaches
95\% SOC during the morning (Fig.~\ref{fig:ntpll_daily_audit}(e)). It then
delivers 219.69 kWh during 10:00--16:00, while EVs receive 138.47 kWh in
that window (panel (f)). Only 46.90 kWh of BESS discharge occurs during the
16:00--21:00 on-peak window. Hence dispatch is not simply ``wait for the
highest TOU price'': the price-weighted wholesale obligation, shared capability,
and contemporaneous meter constraints also affect when stored energy is used.
On July 15, the same site instead takes in 72.76 kWh during 10:00--16:00
and delivers 240.24 kWh during 16:00--21:00
(Fig.~\ref{fig:ntpll_daily_audit_later}(e)). The contrast shows why a single
day is insufficient to characterize the dispatch policy.

At Center Hall, 178.91 of 247.74 kWh of daily EV service is delivered during
10:00--16:00 on July 1, accompanied by 199.10 kWh of BESS discharge
(Fig.~\ref{fig:center_daily_audit}). On July 15, 215.80 of 238.17 kWh of
EV service lies in the same window, but BESS discharge is split more evenly
between 10:00--16:00 (134.59 kWh) and 16:00--21:00 (148.49 kWh)
(Fig.~\ref{fig:center_daily_audit_later}). These quantities are integrated
from executed 15-min power, not inferred from a forecast plan. The observed
coincidence of EV service, BESS response, and price signals explains the
physical mechanism; it does not identify a unique causal contribution of each
objective term without a separate controlled ablation.

\begin{figure*}[!t]
\centering
\includegraphics[width=\textwidth]{fig_08_ntpll_perfect_daily_6panel_20250701_paper.png}
\caption{NTPLL, July 1, 2025: full building--PV--EV--BESS portfolio with the 24-h Perfect EV updater and quarterly CAISO RU/RD attenuation factors. BESS power is positive for charging. Panel (e) retains the WM/NWM charging and discharging channels and the SOC legend; panel (f) retains the original power-variable and threshold labels.}
\label{fig:ntpll_daily_audit}
\end{figure*}
\begin{figure*}[!t]
\centering
\includegraphics[width=\textwidth]{fig_09_ntpll_perfect_daily_6panel_20250715_paper.png}
\caption{NTPLL, July 15, 2025, with the same configuration and corresponding axis limits as Fig.~\ref{fig:ntpll_daily_audit}. Midday replenishment is followed by substantially more evening discharge than on July 1.}
\label{fig:ntpll_daily_audit_later}
\end{figure*}
\begin{figure*}[!t]
\centering
\includegraphics[width=\textwidth]{fig_10_center_hall_perfect_daily_6panel_20250701_paper.png}
\caption{Center Hall, July 1, 2025: the full portfolio, 24-h Perfect updater, and the same attenuation-factor assumptions as NTPLL. Panel definitions follow Fig.~\ref{fig:ntpll_daily_audit}.}
\label{fig:center_daily_audit}
\end{figure*}
\begin{figure*}[!t]
\centering
\includegraphics[width=\textwidth]{fig_11_center_hall_perfect_daily_6panel_20250715_paper.png}
\caption{Center Hall, July 15, 2025, with the same corresponding axis limits as Fig.~\ref{fig:center_daily_audit}. The BESS serves both the midday charging period and the evening on-peak window.}
\label{fig:center_daily_audit_later}
\end{figure*}

\subsection{One-week sensitivity analysis}
\label{sec:attenuation_scope}

The four controlled experiments use NTPLL data for March 1--7, 2025,
including both weekdays and a weekend. Each setting is evaluated with the
Persistence and Perfect updaters in the same 24-h rolling controller. The
reference has a 250-kW/332-kWh BESS, full EV service, retail--energy--AS
participation, measured passive building/PV profiles, and SP/NSP attenuation
factors of 0.2. RU/RD retain the official CAISO hourly attenuation values
\cite{caisoAttenuation2025}, repeated over the four contained 15-min intervals;
these are expected-deployment proxies rather than realized AGC trajectories.
All cases start from the same February Perfect-dispatch history and 50\%
BESS SOC, and end at 50\% SOC. Initial March PD/NCD thresholds are zero.
Subsequent EV baselines and billing states evolve with each executed
trajectory. DA and RT decisions are reoptimized within the existing market
timeline, with submitted DA awards fixed in the corresponding RT problems.
Thus these are closed-loop scenario comparisons, not fixed-DA RT attribution
tests. Site profiles and prices are identical across the paired cases.

Tables~\ref{tab:weekly_factors}--\ref{tab:weekly_markets} report weekly USD:
Net is EV charging-service revenue plus WM revenue and the signed TOU, PD,
and NCD cost entries; WM includes energy settlement and AS capacity payments.
Negative Net includes the passive building electricity bill and should not
be read as standalone EV--BESS profitability. Demand charges are increments
from the initial March peaks, not prorated weekly charges or a completed
monthly bill. Gain is the improvement relative to the same updater's
reference, except in Table~\ref{tab:weekly_markets}, where it is relative to
retail-only operation. Independently rounded daily accounts can produce
cent-level differences between displayed component sums and weekly totals.

\paragraph{SP/NSP attenuation factors}
Constant factors of 0, 0.2, 0.5, and 1 are compared with time-varying synthetic
profiles. Three random permutations span $[0,0.4]$ with an exact daily mean
of 0.2 (seeds 11, 22, and 33); another spans $[0,1]$ with daily mean 0.5
(seed 44). SP and NSP share the same synthetic profile, and each paired
updater uses the identical realization. These profiles are known scenario
inputs, not measurements or tests of attenuation-forecast error. RU/RD are
not varied.

\begin{table*}[p]
\centering
\caption{SP/NSP attenuation sensitivity at NTPLL, March 1--7, 2025. All financial entries are USD. Random settings are labeled by daily mean and seed; RU/RD use the same measured factors in every row.}
\label{tab:weekly_factors}
\footnotesize
\setlength{\tabcolsep}{3pt}
\begin{tabular}{llrrrrrrr}
\toprule
Setting & Updater & Net & WM & TOU & PD & NCD & EV & Gain \\
\midrule
Constant 0.2 (ref.) & Persistence & -47194.06 & 935.99 & -31743.03 & -678.93 & -16574.34 & 866.23 & 0.00 \\
Constant 0 & Persistence & -47177.03 & 941.04 & -31740.15 & -678.56 & -16565.61 & 866.23 & 17.03 \\
Constant 0.5 & Persistence & -47188.94 & 939.46 & -31755.35 & -678.37 & -16560.92 & 866.23 & 5.12 \\
Constant 1 & Persistence & -47190.57 & 939.22 & -31742.72 & -678.92 & -16574.35 & 866.23 & 3.49 \\
Random 0.2; seed 11 & Persistence & -47179.60 & 937.19 & -31741.11 & -678.48 & -16563.45 & 866.23 & 14.46 \\
Random 0.2; seed 22 & Persistence & -47187.03 & 937.55 & -31745.59 & -678.60 & -16566.64 & 866.23 & 7.03 \\
Random 0.2; seed 33 & Persistence & -47179.66 & 937.47 & -31741.56 & -678.47 & -16563.33 & 866.23 & 14.40 \\
Random 0.5; seed 44 & Persistence & -47183.21 & 941.07 & -31758.41 & -678.10 & -16554.03 & 866.23 & 10.85 \\
\midrule
Constant 0.2 (ref.) & Perfect & -47060.28 & 924.66 & -31677.25 & -675.81 & -16498.10 & 866.23 & 0.00 \\
Constant 0 & Perfect & -47091.67 & 921.23 & -31678.57 & -676.86 & -16523.73 & 866.23 & -31.39 \\
Constant 0.5 & Perfect & -47054.99 & 925.52 & -31677.74 & -675.50 & -16493.54 & 866.23 & 5.29 \\
Constant 1 & Perfect & -47059.31 & 926.27 & -31679.88 & -675.51 & -16496.44 & 866.23 & 0.97 \\
Random 0.2; seed 11 & Perfect & -47098.87 & 922.00 & -31681.76 & -677.05 & -16528.30 & 866.23 & -38.59 \\
Random 0.2; seed 22 & Perfect & -47060.21 & 919.78 & -31681.57 & -675.45 & -16489.20 & 866.23 & 0.07 \\
Random 0.2; seed 33 & Perfect & -47060.20 & 922.19 & -31678.66 & -675.66 & -16494.30 & 866.23 & 0.08 \\
Random 0.5; seed 44 & Perfect & -47054.71 & 926.89 & -31678.31 & -675.35 & -16494.19 & 866.23 & 5.57 \\
\bottomrule
\end{tabular}
\end{table*}

Table~\ref{tab:weekly_factors} shows modest changes in total operating value:
relative to the 0.2 reference, Net changes by $+3.49$ to $+17.03$ USD for
Persistence and $-38.59$ to $+5.57$ USD for Perfect. This does not imply
that attenuation is financially irrelevant. Combining
\eqref{eq:wm_energy} and \eqref{eq:wm_deployment_revenue}, the interval
energy receipt can be written as
\begin{align}
r_t^{\WM,\mathrm{energy}}
=\dt\left[(\lambda_t^{\DA}-\lambda_t^{\RT})p_t^{\DA}
+\lambda_t^{\RT}q_t^{\WM}\right].
\label{eq:sensitivity_energy_identity}
\end{align}
Here $r_t^{\WM,\mathrm{energy}}$ is the interval energy revenue in USD and
$q_t^{\WM}$ is the attenuation-adjusted wholesale obligation in kW defined
in \eqref{eq:qwm}. At a fixed DA energy position and fixed obligation, an
attenuation change can be offset by a signed RT energy adjustment, leaving
this energy receipt unchanged. Capacity payments, shared power limits, SOC,
and the retail bill can nevertheless change when the controller reoptimizes.
The observed small net changes are therefore consistent with substitution
between energy and reserve positions; they are not a general invariance
result. With the retained 1\% MILP tolerance, these small differences also
do not establish a reliable ranking of factor settings.

\paragraph{BESS power and energy ratings}
\begin{table*}[p]
\centering
\caption{BESS rating sensitivity at NTPLL, March 1--7, 2025 (USD). Gain is relative to the 250-kW/332-kWh reference for the corresponding updater. Investment costs are excluded.}
\label{tab:weekly_bess}
\footnotesize
\setlength{\tabcolsep}{3pt}
\begin{tabular}{llrrrrrrr}
\toprule
Setting & Updater & Net & WM & TOU & PD & NCD & EV & Gain \\
\midrule
250 kW / 332 kWh (ref.) & Persistence & -47194.06 & 935.99 & -31743.03 & -678.93 & -16574.34 & 866.23 & 0.00 \\
500 kW / 332 kWh & Persistence & -46483.55 & 1689.70 & -31727.58 & -681.23 & -16630.68 & 866.23 & 710.51 \\
250 kW / 664 kWh & Persistence & -46093.02 & 1015.29 & -31214.41 & -659.13 & -16101.00 & 866.23 & 1101.04 \\
500 kW / 664 kWh & Persistence & -45329.62 & 1833.31 & -31222.48 & -661.35 & -16145.32 & 866.23 & 1864.44 \\
\midrule
250 kW / 332 kWh (ref.) & Perfect & -47060.28 & 924.66 & -31677.25 & -675.81 & -16498.10 & 866.23 & 0.00 \\
500 kW / 332 kWh & Perfect & -46296.62 & 1690.01 & -31684.96 & -675.58 & -16492.34 & 866.23 & 763.66 \\
250 kW / 664 kWh & Perfect & -45727.96 & 1013.72 & -31168.64 & -646.90 & -15792.40 & 866.23 & 1332.32 \\
500 kW / 664 kWh & Perfect & -44944.71 & 1818.23 & -31182.93 & -646.89 & -15799.36 & 866.23 & 2115.57 \\
\bottomrule
\end{tabular}
\end{table*}

Doubling only power raises WM receipts from 935.99 to 1689.70 USD under
Persistence and from 924.66 to 1690.01 USD under Perfect
(Table~\ref{tab:weekly_bess}). Doubling only energy gives much smaller WM
increases, to 1015.29 and 1013.72 USD, respectively, but larger reductions in
TOU and NCD costs. The distinction is operationally meaningful: power rating
expands instantaneous market capability, whereas energy capacity supports
longer load shifting and peak management. Doubling both improves weekly Net
by 1864.44 USD for Persistence and 2115.57 USD for Perfect. These results
support studying power and energy separately rather than using a single
battery-size multiplier. They are operating-value comparisons, not optimal
investment recommendations, because additional capital and financing costs
are not included.

\paragraph{Minimum EV service}
Only in this experiment, the full-service equality in \eqref{eq:ev_core}
is replaced by
\begin{align}
\eta_{\mathrm{srv}}E_i^{\req}
\le E_i^{\mathrm{del}} \le E_i^{\req},\qquad
\eta_{\mathrm{srv}}\in\{1,0.8,0.5\},
\label{eq:sensitivity_service}
\end{align}
where $E_i^{\mathrm{del}}$ is delivered session energy in kWh, calculated
with the same convention as \eqref{eq:ev_core}, and
$\eta_{\mathrm{srv}}$ is the minimum service fraction, distinct from the
charging efficiency $\eta_i$. A lower minimum permits, but does not require,
less charging. All other experiments retain full service. Driver payments
are calculated from actual delivered charging energy, not requested energy.
\begin{table*}[p]
\centering
\caption{Minimum-service sensitivity at NTPLL, March 1--7, 2025 (USD). EV revenue is based on delivered charging energy; gains do not include compensation for unserved demand.}
\label{tab:weekly_eta}
\footnotesize
\setlength{\tabcolsep}{3pt}
\begin{tabular}{llrrrrrrr}
\toprule
Setting & Updater & Net & WM & TOU & PD & NCD & EV & Gain \\
\midrule
$\eta_{\mathrm{srv}}=1$ (ref.) & Persistence & -47194.06 & 935.99 & -31743.03 & -678.93 & -16574.34 & 866.23 & 0.00 \\
$\eta_{\mathrm{srv}}=0.8$ & Persistence & -47057.29 & 939.80 & -31616.15 & -673.03 & -16430.27 & 722.35 & 136.77 \\
$\eta_{\mathrm{srv}}=0.5$ & Persistence & -46910.94 & 941.95 & -31484.87 & -664.87 & -16231.46 & 528.31 & 283.12 \\
\midrule
$\eta_{\mathrm{srv}}=1$ (ref.) & Perfect & -47060.28 & 924.66 & -31677.25 & -675.81 & -16498.10 & 866.23 & 0.00 \\
$\eta_{\mathrm{srv}}=0.8$ & Perfect & -46903.49 & 932.85 & -31565.20 & -669.74 & -16352.82 & 751.41 & 156.79 \\
$\eta_{\mathrm{srv}}=0.5$ & Perfect & -46761.54 & 932.95 & -31479.49 & -662.46 & -16172.20 & 619.63 & 298.74 \\
\bottomrule
\end{tabular}
\end{table*}

At minimum service 0.8, delivered energy is 83.39\% of requested energy for
Persistence and 86.74\% for Perfect; at 0.5, these fractions are 60.99\% and
71.53\%, respectively. Thus the optimizer does not simply enforce equality
at the minimum. Table~\ref{tab:weekly_eta} shows Net improvements of 136.77
and 156.79 USD at 0.8, increasing to 283.12 and 298.74 USD at 0.5, for
Persistence and Perfect, respectively. The latter gains accompany 844.82
and 616.52 kWh of unserved demand and lower driver payments. They must
therefore be interpreted as a service--cost trade-off, not a free efficiency
gain. Customer acceptance and any compensation would be needed before
adopting such a policy.

\paragraph{Joint retail--wholesale participation}
\begin{table*}[p]
\centering
\caption{Market-participation comparison at NTPLL, March 1--7, 2025 (USD). Gain uses the corresponding retail-only case, not the full-market reference.}
\label{tab:weekly_markets}
\footnotesize
\setlength{\tabcolsep}{3pt}
\begin{tabular}{llrrrrrrr}
\toprule
Setting & Updater & Net & WM & TOU & PD & NCD & EV & Gain \\
\midrule
Retail only & Persistence & -48145.54 & 0.00 & -31725.41 & -680.22 & -16606.13 & 866.23 & 0.00 \\
Retail + energy & Persistence & -47613.71 & 512.66 & -31729.63 & -679.31 & -16583.69 & 866.23 & 531.83 \\
Retail + energy + AS & Persistence & -47194.06 & 935.99 & -31743.03 & -678.93 & -16574.34 & 866.23 & 951.48 \\
\midrule
Retail only & Perfect & -47947.99 & 0.00 & -31672.47 & -674.54 & -16467.21 & 866.23 & 0.00 \\
Retail + energy & Perfect & -47475.82 & 505.87 & -31674.14 & -675.81 & -16497.97 & 866.23 & 472.17 \\
Retail + energy + AS & Perfect & -47060.28 & 924.66 & -31677.25 & -675.81 & -16498.10 & 866.23 & 887.71 \\
\bottomrule
\end{tabular}
\end{table*}

Adding energy trading to retail operation improves Net by 531.83 USD under
Persistence and 472.17 USD under Perfect; adding AS increases the total
gains to 951.48 and 887.71 USD (Table~\ref{tab:weekly_markets}). Hence AS
adds 419.65 and 415.54 USD beyond the energy-only controller in this week.
The gain is not equal to gross WM receipts because market actions also
change TOU costs and monthly peaks. For example, under Persistence the
energy-only case earns 512.66 USD in WM receipts and improves Net by
531.83 USD after retail interactions. The comparison quantifies the value
of joint optimization, rather than treating wholesale income as an
independent revenue stream that can simply be added to an unchanged bill.

\paragraph{Separate full-horizon benchmark}
Table~\ref{tab:weekly_horizon} uses a separate, common pre-period-derived EV
baseline frozen throughout the week, so its 24-h comparator values should
not be substituted for the evolving-baseline references above. Site inputs,
initial billing thresholds, full EV service, BESS rating, and initial/final
SOC are otherwise matched. The oracle sees the entire week in one problem.
\begin{table*}[p]
\centering
\caption{Separate common-baseline horizon comparison for NTPLL, March 1--7, 2025 (USD). Oracle revenue is a feasible incumbent, not its solver bound.}
\label{tab:weekly_horizon}
\footnotesize
\setlength{\tabcolsep}{5pt}
\begin{tabular}{lrrrrrr}
\toprule
Method & Net & WM & TOU & PD & NCD & EV \\
\midrule
24-h Persistence & -47185.99 & 930.84 & -31744.07 & -678.36 & -16560.64 & 866.23 \\
24-h Perfect & -47055.79 & 922.88 & -31682.77 & -675.01 & -16487.15 & 866.23 \\
Full-horizon oracle & -46842.05 & 945.95 & -31514.96 & -674.44 & -16464.84 & 866.23 \\
\bottomrule
\end{tabular}
\end{table*}

The oracle incumbent exceeds these executed Persistence and Perfect revenues
by 343.94 and 213.74 USD, respectively. Its remaining maximization bound is
$-46720.39$ USD, versus the incumbent $-46842.05$ USD, corresponding to a
0.260\% gap. There is an additional feasible-set qualification: the oracle
uses true-availability DA offer bounds, while forecast-driven Persistence DA
awards exceed that benchmark's downward offer envelope by up to 73.216 kW,
even though the executed RT dispatch remains physically feasible. This
oracle is therefore reported as a full-information benchmark, not a
universal upper bound on all forecast-driven DA schedules. The common-set
upper-bound argument discussed earlier applies only when its feasible-set
conditions hold.

Across the 32 saved MPC cases underlying these comparisons, all 21728 DA/RT
optimizations met the requested 1\% gap without a time-limit termination.
Saved checks cover service, meter balance, SOC, DA award granularity,
settlement, and boundary conditions. These checks support the internal
consistency of this one-week experiment; they do not establish statistical
robustness across seasons or remove the limitations of expected AS deployment.

\clearpage
\subsection{Numerical acceptance}
\label{sec:site_validation_results}

The 20 monthly cases contain all 58,752 expected RT records and pass the
saved solver, meter, settlement, service, and boundary-state checks. All saved actual AS quantities satisfy nonnegativity within the implemented $10^{-5}$-kW numerical tolerance; that tolerance is not a negative-capacity modeling allowance. The largest
saved meter-balance residual is below $10^{-6}$ kW, and the daily rounded
financial residual remains within the \$0.02 criterion. These checks establish
numerical and accounting consistency of these trajectories, not exact
global optimality or compliance with omitted market qualification rules.

\section{Discussion}
\label{sec:discussion}

\subsection{Site scale changes the value stack}

The two sites demonstrate why a campus-wide total is not an adequate unit of
inference. NTPLL's passive building bill dominates its absolute net value, and
the common BESS is modest relative to the building peak. Center Hall has a much
smaller native net-load peak and a larger PV-to-load ratio, so the same BESS and
a similarly sized charger population have greater leverage over the PCC. The
appropriate comparison is therefore incremental value within each site, not
the absolute dollar level across sites.

The staged calculation also separates physical additions from optimization.
PV changes retail energy and demand charges without receiving a wholesale bid.
The EV--BESS stage then adds flexible dispatch, EV-service revenue, and WM
participation. This ordering avoids crediting the controller with passive PV
production and avoids crediting PV or building demand as callable AS capacity.

\subsection{Forecast value is a closed-loop quantity}

The site- and month-dependent forecast rankings in Table~\ref{tab:site_monthly_results} are informative rather than contradictory. A 24-h Perfect updater removes EV forecast error inside the
active window, but the controller still minimizes a truncated objective. At a
site where retail peaks, BESS SOC, and DA commitments interact strongly, the
first action selected by a locally better-informed problem can leave a worse
out-of-window state. Consequently, prediction accuracy and billing-period
economic value must both be reported. A future forecast model should be judged
by executed closed-loop value under matched physical and financial checks, not
by forecast error alone.

\subsection{Operational interpretation}

The results support three bounded application conclusions. First, WM and
retail accounts must be assessed jointly: the June Center Hall decomposition
in Section~\ref{sec:site_forecast_results} shows why demand-charge relief alone
does not determine net value. Second, the separate DA/RT bars and combined
obligation in Figs.~\ref{fig:ntpll_daily_audit}(c)--(d) and
\ref{fig:center_daily_audit}(c)--(d) distinguish financial recourse from actual
resource delivery; the shared-capability equations specify feasibility rather
than guaranteeing eligibility for a real market. Third, the power scales in
Fig.~\ref{fig:site_power_specs} and the component accounts in
Fig.~\ref{fig:site_resource_financial} show why the same BESS has different
relative influence at the two sites. The July 1/15 dispatch contrasts further
demonstrate that the allocation of stored energy between midday and evening
is state dependent. These observations motivate site-specific evaluation, not
a universal forecast ranking or storage-sizing recommendation.

\section{Limitations and future work}
\label{sec:limitations}

The formal site experiment covers two UCSD configurations and five consecutive
billing months spanning the modeled winter and summer schedules, not a full year. The same 250-kW/332-kWh BESS is imposed on both sites to isolate
data-scale effects; it is not an economically optimized site-specific size.
The PV capacity values are observed proxies, not certified nameplate ratings.
The building channels are native net meters, so gross demand is reconstructed
using the mapped PV signal rather than measured by a separate gross-load meter.
Feeder topology, transformer limits, detailed BESS degradation, and V2G export
are outside the present model.

Market prices are treated as known over each active horizon. Expected AS
deployment replaces realized AGC, and the model omits regulation mileage,
performance scores, no-pay charges, telemetry, minimum bid sizes, and
scheduling-coordinator qualification. The site study does not recompute a full
billing-period oracle; Perfect information remains a 24-h EV-information
benchmark. Future work should add further sites and seasons, price forecasts,
realized deployment data, and site-specific BESS sizing while retaining the
same meter, settlement, service, and terminal-state validation gates.

\section{Conclusion}
\label{sec:conclusion}

The contribution of this work is an integrated, application-oriented
optimization framework rather than an additional resource modeled in isolation.
A DA optimization commits hourly energy and AS awards, and successive RT
optimizations update EV information and choose 15-min recourse while honoring
those commitments. EV charging, BESS power and energy, passive PV/building
net load, retail TOU and demand charges, and wholesale energy and capacity
therefore interact within one physically constrained value stack. Relative to
EV-focused predictive scheduling \cite{chen2026realtime}, the framework brings
storage and the site meter into the same DA--RT economic decisions; relative
to BTM storage revenue stacking \cite{seward2022revenueStacking}, it adds
session-level charging requirements and dynamically updated EV availability.
The contribution is the joint treatment of these couplings, not complexity
for its own sake or a claim that each individual modeling component is new.

The measured March--July results demonstrate why this integration matters.
Against the same site's building-plus-PV counterfactual, the deployable
Persistence portfolio improves five-month net value by \$56,873.35 at NTPLL
and \$43,149.11 at Center Hall. These improvements comprise \$18,909.61
and \$19,893.52 of WM revenue, \$13,012.41 and \$17,419.60 of EV-service
receipts, and \$24,951.34 and \$5,835.99 of retail-cost reductions,
respectively (Table~\ref{tab:site_monthly_decomposition}). Thus, even after
adding EV charging demand, the full portfolio reduces the modeled retail bill
by 2.72\% and 8.42\% over the five months. These are operating-account
comparisons, not isolated controller savings against an otherwise identical
unmanaged EV--BESS portfolio or capital-investment returns.

The cross-month and daily decompositions extend the tariff-dependent
PV--storage perspective \cite{boampong2020retail}: the same BESS rating leads
to different retail/wholesale trade-offs at different sites, and stored energy
can support midday EV service on one day and evening tariff relief on another.
Perfect 24-h EV information does not uniformly improve billing-period value,
because the forecast is embedded in a stateful, finite-horizon controller.
Reporting only forecast accuracy or one summer-month total would miss these
application-relevant effects. Comparisons with prior studies are therefore
based on modeling scope and observed mechanisms, not incomparable revenue
percentages under different tariffs and portfolios.

For an aggregator, the results quantify multiple operating value streams while
preserving charging service. For the grid interface, the formulation makes
energy and reserve commitments share an explicit deliverability envelope,
providing a route to coordinated flexibility rather than independent,
potentially conflicting bids. Network-level reliability, avoided emissions,
and realized AGC performance are not measured here, so these prospective
benefits are not claimed as demonstrated outcomes. Within these limits, the
two-site, five-month results support a practical framework for evaluating
retail--wholesale coordination across heterogeneous BTM portfolios.

\appendix

\section{Complete optimization model}
\label{app:complete_model}

This appendix gives the index-complete formulation corresponding to the compact
coupling equations in Section~\ref{sec:model}. It also distinguishes physical
variables from accounting variables, which is essential when an EV load
reduction is labeled ``discharge'' relative to baseline.

\subsection{EV feasibility and accounting channels}
Let $y_{t,i}$ be energy delivered to session $i$ in interval $t$, and let
$\bar y_{t,i}$ be the availability-dependent interval limit. The implementation
uses
\begin{align}
0&\le y_{t,i}\le \bar y_{t,i}, &&t\in\calT,\ i\in\calI, \label{eq:app_ev_interval}\\
\sum_{i\in\calI}y_{t,i}&=p_t^{\EV}\dt, &&t\in\calT, \label{eq:app_ev_aggregate}\\
\sum_{t\in\calT_i}y_{t,i}&=E_i^{\req}, &&i\in\calI. \label{eq:app_ev_energy}
\end{align}

Equation~\eqref{eq:app_ev_interval} limits each session to its connected interval capability; both delivered and available interval energies are in kWh, so an absent session has zero capability; equation~\eqref{eq:app_ev_aggregate} converts aggregate delivered interval energy into physical EV charging power in kW through the interval duration in hours; equation~\eqref{eq:app_ev_energy} enforces the requested service by departure. In an RT subproblem, already executed energy is deducted from the remaining requirement, rather than being delivered again.
Therefore every compared policy delivers the same requested session energy
when its EV-service checks pass. Relative to baseline $B_t^{\EV}$, the EV
deviation is partitioned into wholesale and non-wholesale channels:
\begin{align}
p_t^{\EV}-B_t^{\EV}
={}&p_t^{\ch,\EV,\WM}-p_t^{\dch,\EV,\WM}\nonumber\\
&+p_t^{\ch,\EV,\NWM}-p_t^{\dch,\EV,\NWM}, \label{eq:app_ev_split}\\
p_t^{\dch,\EV,\WM}+p_t^{\dch,\EV,\NWM}
&\le M_t^{\EV}u_t^{\dch,\EV}, \label{eq:app_ev_dch_binary}\\
p_t^{\ch,\EV,\WM}+p_t^{\ch,\EV,\NWM}
&\le M_t^{\EV}u_t^{\ch,\EV}, \label{eq:app_ev_ch_binary}\\
u_t^{\ch,\EV}+u_t^{\dch,\EV}&\le1. \label{eq:app_ev_direction}
\end{align}

Equation~\eqref{eq:app_ev_split} partitions the signed deviation from the EV reference profile into WM and NWM accounting channels without creating physical vehicle discharge; equation~\eqref{eq:app_ev_dch_binary} allows baseline-relative load reduction only when the reduction-direction binary is enabled; the physical bound in kW limits the sum of both accounting channels; equation~\eqref{eq:app_ev_ch_binary} similarly gates charging above baseline across the two channels; equation~\eqref{eq:app_ev_direction} prevents simultaneous positive and negative baseline deviations in the same interval; the direction indicators are binary.
All four channel variables are nonnegative. The ``discharging'' term in
\eqref{eq:app_ev_split} is load reduction only; the physical EV power in
\eqref{eq:app_ev_aggregate} remains nonnegative.

\subsection{BESS physics, channel split, and terminal state}
The BESS variables obey
\begin{align}
p_t^{\BESS}&=p_t^{\ch,\BESS}-p_t^{\dch,\BESS}, \label{eq:app_bess_net}\\
p_t^{\ch,\BESS}&=p_t^{\ch,\BESS,\WM}
+p_t^{\ch,\BESS,\NWM}, \label{eq:app_bess_ch_split}\\
p_t^{\dch,\BESS}&=p_t^{\dch,\BESS,\WM}
+p_t^{\dch,\BESS,\NWM}, \label{eq:app_bess_dch_split}\\
0\le p_t^{\ch,\BESS}&\le P^{\BESS}u_t^{\ch,\BESS}, \label{eq:app_bess_charge_power}\\
0\le p_t^{\dch,\BESS}&\le P^{\BESS}u_t^{\dch,\BESS}, \label{eq:app_bess_discharge_power}\\
u_t^{\ch,\BESS}+u_t^{\dch,\BESS}&\le1, \label{eq:app_bess_direction}\\
s_{t+1}&=s_t+\frac{\dt\sqrt{\gamma}}{C^{\BESS}}
p_t^{\ch,\BESS}\nonumber\\
&\quad-\frac{\dt}{C^{\BESS}\sqrt{\gamma}}
p_t^{\dch,\BESS}, \label{eq:app_bess_soc}\\
s^{\min}\le s_t&\le s^{\max}. \label{eq:app_bess_soc_bounds}
\end{align}

Equation~\eqref{eq:app_bess_net} defines positive BESS net power as charging and negative power as discharging; equation~\eqref{eq:app_bess_ch_split} allocates total physical charging between WM and NWM without duplicating energy; equation~\eqref{eq:app_bess_dch_split} does the same for physical discharging; all four channel powers are nonnegative; equation~\eqref{eq:app_bess_charge_power} limits enabled charging by the common inverter power rating in kW; equation~\eqref{eq:app_bess_discharge_power} limits enabled discharging by that same rating; equation~\eqref{eq:app_bess_direction} excludes simultaneous physical charging and discharging; equation~\eqref{eq:app_bess_soc} updates the dimensionless energy state: charging power times hours and charging efficiency adds stored kWh, while discharge removes kWh divided by discharge efficiency; division by rated kWh yields the SOC change; equation~\eqref{eq:app_bess_soc_bounds} reserves the minimum and maximum permitted stored-energy margins in every state.
For each calendar day $d$ intersecting the active horizon,
\begin{align}
\dt\sum_{t\in\calT_d}
\left(p_t^{\ch,\BESS}+p_t^{\dch,\BESS}\right)
&\le2C^{\BESS}(s^{\max}-s^{\min}). \label{eq:app_bess_throughput}
\end{align}

Equation~\eqref{eq:app_bess_throughput} limits daily charge-plus-discharge energy to one usable-capacity cycle, i.e., twice the usable kWh range; any executed throughput belonging to the constrained day must be included when applying the remaining-day budget.
The executed SOC from one call becomes the next call's initial state. At the final billing-period boundary, the corresponding horizon enforces
\begin{align}
s_{T+1}=s^{\mathrm{terminal}}. \label{eq:app_terminal_soc}
\end{align}

Equation~\eqref{eq:app_terminal_soc} fixes the billing-period boundary energy. It does not reset SOC after every day; within a month, executed SOC is passed continuously to the next optimization.

\subsection{Signed meter and passive PV/building sensitivity}
Introduce a configuration indicator $\zeta^{\mathrm{BTM}}\in\{0,1\}$. Both
the reference and sensitivity cases use the same balance
\begin{align}
p_t^{\GI}=p_t^{\EV}+p_t^{\BESS}
+\zeta^{\mathrm{BTM}}(p_t^{\load}-p_t^{\PV}). \label{eq:app_meter_switch}
\end{align}

Equation~\eqref{eq:app_meter_switch} adds the passive net load to physical EV and BESS power. The configuration flag is an input, not a decision variable.
The site-specific results use $\zeta^{\mathrm{BTM}}=1$; setting it to zero
recovers the EV--BESS reference. Native building-meter readings retain signed
net load. The gross-load reconstruction in \eqref{eq:site_net_meter_identity}
prevents double subtraction of PV. Input quality control aligns the channels
to the model's 96-interval local-day clock and flags imputed timestamps,
including daylight-saving gaps; this is a local-clock approximation rather
than an elapsed-time model of a 23-hour day. PV measurement noise is treated
separately from valid negative native net-meter readings. No passive-resource
market baseline is constructed.

The passive-resource market scope is imposed explicitly:
\begin{align}
p_t^{\load,\WM}=p_t^{\PV,\WM}&=0,\label{eq:app_passive_wm_energy_zero}\\
c_t^{x,\load}=c_t^{x,\PV}&=0,
\qquad x\in\calK,\label{eq:app_passive_as_zero}\\
R^{\load,\WM}=R^{\PV,\WM}&=0.\label{eq:app_passive_wm_revenue_zero}
\end{align}

Equation~\eqref{eq:app_passive_wm_energy_zero} excludes passive building/PV from wholesale energy positions; equation~\eqref{eq:app_passive_as_zero} excludes them from all callable AS products; equation~\eqref{eq:app_passive_wm_revenue_zero} therefore assigns no direct WM receipts to these passive resources; their indirect influence on optimized EV--BESS bids remains possible.
These are not additional optimization variables in the implementation; the
equalities state their exclusion from the model. Consequently the EV--BESS
position $p_t^{\actual}=p_t^{\DA}+p_t^{\RT}$ and its physical limits are
identical in the reference and passive-meter cases. The switch changes only
$p_t^{\GI}$, signed TOU settlement, and the PD/NCD threshold trajectory.

\subsection{Native award granularity on the 15-min model clock}
Let $H(t)$ map interval $t$ to its containing DA trade hour. The exact block
constraints used for submitted DA quantities are
\begin{align}
p_t^{\DA}-p_{t-1}^{\DA}&=0,
&&H(t)=H(t-1),\label{eq:app_da_energy_hourly_blocks}\\
c_t^{x,\DA}-c_{t-1}^{x,\DA}&=0,
&&H(t)=H(t-1),\quad x\in\calK.
\label{eq:app_da_hourly_blocks}
\end{align}

Equation~\eqref{eq:app_da_energy_hourly_blocks} holds submitted DA energy power constant within its trade hour; four quarter-hour energy deliveries form one hourly award; equation~\eqref{eq:app_da_hourly_blocks} imposes the same hourly consistency on each DA capacity product, while RT recourse remains interval-specific.
No analogous equality is applied to RT adjustments. If native energy prices
$\widetilde\lambda$ are US\$/MWh and native capacity prices
$\widetilde\rho$ are US\$/MW per award interval, the arrays passed to the
optimization are
\begin{align}
\lambda_t^{s}&=\frac{\widetilde\lambda_t^{s}}{10^3\ \mathrm{kWh/MWh}},
&&s\in\{\DA,\RT\},\label{eq:app_lmp_units}\\
\rho_t^{x,\DA}&=\frac{\widetilde\rho_{H(t)}^{x,\DA}}{(10^3\ \mathrm{kW/MW})(1\ \mathrm{h})},
&&\label{eq:app_da_as_units}\\
\rho_t^{x,\RT}&=\frac{\widetilde\rho_t^{x,\RT}}{(10^3\ \mathrm{kW/MW})\dt}.
&&\label{eq:app_rt_as_units}
\end{align}

Equation~\eqref{eq:app_lmp_units} converts energy prices from dollars per MWh to dollars per kWh without changing the energy-delivery duration; equation~\eqref{eq:app_da_as_units} spreads the hourly DA capacity award over its one-hour duration. Repeating this coefficient four times and multiplying each by 0.25 h pays the hourly award once; equation~\eqref{eq:app_rt_as_units} normalizes each native RT capacity award by its own 0.25-h duration; multiplying by the model interval duration recovers that single binding award payment.
Every displayed price table and panel uses these normalized coefficients. The
single settlement expression $\dt\rho c$ then reproduces both a one-hour DA
award and one 15-min RT binding award, as proved in
\eqref{eq:da_capacity_payment_check}--\eqref{eq:rt_capacity_payment_check}.

\subsection{DA/RT settlement and shared deliverability}
For each interval, the cleared DA arrays already contain zero whenever no award
is active. They are fixed directly in RT:
\begin{align}
p_t^{\DA}&=\bar p_t^{\DA}, \label{eq:app_fixed_da_energy}\\
c_t^{x,\DA}&=\bar c_t^{x,\DA}, \label{eq:app_fixed_da_as}\\
p_t^{\actual}&=p_t^{\DA}+p_t^{\RT}, \label{eq:app_actual_energy}\\
c_t^{x,\actual}&=c_t^{x,\DA}+c_t^{x,\RT}\ge0. \label{eq:app_actual_as}
\end{align}

Equation~\eqref{eq:app_fixed_da_energy} makes already cleared DA energy a fixed parameter in RT, preserving market chronology; equation~\eqref{eq:app_fixed_da_as} similarly preserves each cleared DA AS award instead of retrospectively replacing it; equation~\eqref{eq:app_actual_energy} defines the resulting signed energy position after incremental RT trading; equation~\eqref{eq:app_actual_as} allows a signed RT correction but prohibits negative physically supplied capacity.
The RT variables are signed. In particular, $c_t^{x,\RT}<0$ records a
downward correction to a DA AS award; it does not permit negative actual
capacity. The DA awards and actual positions share the resource envelope:
\begin{align}
c_t^{\RU,\DA}+c_t^{\SP,\DA}+c_t^{\NSP,\DA}
&\le\bar P_t^{\mathrm{up}}, \label{eq:app_da_up_capacity}\\
c_t^{\RD,\DA}&\le\bar P_t^{\mathrm{down}}, \label{eq:app_da_down_capacity}\\
\pos{p_t^{\actual}}+c_t^{\RU,\actual}+c_t^{\SP,\actual}+c_t^{\NSP,\actual}
&\le\bar P_t^{\mathrm{up}}, \label{eq:app_actual_up_capacity}\\
\pos{-p_t^{\actual}}+c_t^{\RD,\actual}
&\le\bar P_t^{\mathrm{down}}, \label{eq:app_actual_down_capacity}\\
-\bar P_t^{\mathrm{down}}\le p_t^{\DA}
&\le\bar P_t^{\mathrm{up}}, \label{eq:app_da_energy_bounds}\\
-\bar P_t^{\mathrm{down}}\le p_t^{\RT}
&\le\bar P_t^{\mathrm{up}}. \label{eq:app_rt_energy_bounds}
\end{align}

Equation~\eqref{eq:app_da_up_capacity} checks the sum of upward DA reserves against upward EV--BESS capability; equation~\eqref{eq:app_da_down_capacity} checks downward DA reserve against downward capability; equation~\eqref{eq:app_actual_up_capacity} allocates upward headroom jointly to positive actual energy and the three upward reserve products; the positive-part operator avoids treating imports as negative consumption of this headroom; equation~\eqref{eq:app_actual_down_capacity} allocates downward headroom jointly to actual imports and RD capacity; equation~\eqref{eq:app_da_energy_bounds} bounds the signed DA energy account by the resource envelope; equation~\eqref{eq:app_rt_energy_bounds} also bounds the signed RT adjustment account; this additional account-level restriction is distinct from the limits on the combined actual position.
Here
\begin{align}
\bar P_t^{\mathrm{up}}&=P^{\BESS}+B_t^{\EV},\label{eq:app_up_envelope}\\
\bar P_t^{\mathrm{down}}&=\max\{P^{\BESS}-B_t^{\EV}+P_{t,\max}^{\EV},0\}.\label{eq:app_down_envelope}
\end{align}


Equation~\eqref{eq:app_up_envelope} adds BESS discharge capability to the EV
baseline reduction available without V2G; equation~\eqref{eq:app_down_envelope}
adds BESS charging capability to EV charging headroom above baseline, with a
nonnegative floor. Both are power bounds in kW; passive load and PV do not
increase these controllable market envelopes.

Panel (c) plots $\bar P_t^{\mathrm{up}}$ and
$-\bar P_t^{\mathrm{down}}$ for both the reference and passive-meter cases.
These are physical EV--BESS capability lines, not graphical envelopes modified
to enclose each separately drawn bid account.

\subsection{Expected deployment and net-flow direction linearization}
The total deployed EV--BESS wholesale obligation is
\begin{align}
q_t^{\WM}={}&p_t^{\DA}+p_t^{\RT}
+\alpha_t^{\RU}c_t^{\RU,\actual}
-\alpha_t^{\RD}c_t^{\RD,\actual}\nonumber\\
&+\alpha_t^{\SP}c_t^{\SP,\actual}
+\alpha_t^{\NSP}c_t^{\NSP,\actual}. \label{eq:app_obligation}
\end{align}

Equation~\eqref{eq:app_obligation} combines the actual energy schedule with expected AS energy expressed as average interval power. RD enters negatively because its expected delivery absorbs rather than exports energy.
Nonnegative variables $p_t^{\ch,\net}$ and $p_t^{\dch,\net}$ select the
negative and positive parts of $q_t^{\WM}$. The exact
implementation is
\begin{align}
u_t^{\ch,\net}+u_t^{\dch,\net}&\le1, \label{eq:app_net_direction}\\
p_t^{\dch,\net}&\ge q_t^{\WM}, \label{eq:app_dch_lower}\\
p_t^{\dch,\net}&\le q_t^{\WM}+\bar P_t^{\mathrm{up}}u_t^{\ch,\net}, \label{eq:app_dch_upper_q}\\
p_t^{\dch,\net}&\le\bar P_t^{\mathrm{up}}u_t^{\dch,\net}, \label{eq:app_dch_upper_binary}\\
p_t^{\ch,\net}&\ge-q_t^{\WM}, \label{eq:app_ch_lower}\\
p_t^{\ch,\net}&\le-q_t^{\WM}+\bar P_t^{\mathrm{down}}u_t^{\dch,\net}, \label{eq:app_ch_upper_q}\\
p_t^{\ch,\net}&\le\bar P_t^{\mathrm{down}}u_t^{\ch,\net}, \label{eq:app_ch_upper_binary}\\
p_t^{\ch,\net}-p_t^{\dch,\net}
&=p_t^{\ch,\BESS,\WM}-p_t^{\dch,\BESS,\WM}\nonumber\\
&\quad+p_t^{\ch,\EV,\WM}-p_t^{\dch,\EV,\WM}. \label{eq:app_net_link}
\end{align}

Equation~\eqref{eq:app_net_direction} selects at most one direction for the net wholesale flow; equation~\eqref{eq:app_dch_lower} makes the nonnegative discharge variable at least the signed obligation; equation~\eqref{eq:app_dch_upper_q} makes that discharge variable no greater than the obligation when charging direction is disabled; equation~\eqref{eq:app_dch_upper_binary} forces net discharge to zero when its binary is off and otherwise limits it by physical upward capability; equation~\eqref{eq:app_ch_lower} makes the nonnegative charge variable at least the negative obligation; equation~\eqref{eq:app_ch_upper_q} makes net charge no greater than the negative obligation when discharge direction is disabled; equation~\eqref{eq:app_ch_upper_binary} forces net charge to zero when its binary is off and otherwise limits it by downward capability; equation~\eqref{eq:app_net_link} equates the net wholesale flow to the sum of EV and BESS WM channels. Together, the direction inequalities yield discharge equal to the positive part of the obligation and charge equal to its negative part; they do not add a tracking penalty.
Equations~\eqref{eq:app_dch_lower}--\eqref{eq:app_ch_upper_binary} impose the
positive/negative-part direction exactly; their coefficients are the physical
capability bounds, not objective penalties on RT movement.

For assumed service duration $T^{\mathrm{DU}}$, the conservative BESS energy
guards are
\begin{align}
s_{t+1}&\ge s^{\min}
+\frac{T^{\mathrm{DU}}}{C^{\BESS}\sqrt{\gamma}}
\left(c_t^{\RU,\actual}+c_t^{\SP,\actual}\right.\nonumber\\
&\hspace{28mm}\left.+c_t^{\NSP,\actual}\right),
\label{eq:app_duration_up}\\
s_{t+1}&\le s^{\max}
-\frac{T^{\mathrm{DU}}\sqrt{\gamma}}{C^{\BESS}}c_t^{\RD,\actual}.
\label{eq:app_duration_down}
\end{align}

Equation~\eqref{eq:app_duration_up} holds enough stored energy above minimum SOC to sustain the assumed upward service duration, including discharge losses; equation~\eqref{eq:app_duration_down} holds enough empty storage below maximum SOC to absorb the assumed downward service energy, including charging efficiency.

\subsection{Retail epigraphs, wholesale revenue, and full objective}
For $q\in\{\PD,\NCD\}$ and eligible interval set $\calT_q$,
\begin{align}
z_k^q&\ge0,&
z_k^q&\ge p_t^{\GI}-M_k^{q,\hist},\quad
t\in\calH_k\cap\calT_q. \label{eq:app_demand_epigraph}
\end{align}

Equation~\eqref{eq:app_demand_epigraph} defines a nonnegative incremental peak in kW above the carried historical threshold. Minimization with a positive demand rate makes it equal to the largest eligible exceedance, or zero if there is none.
The deployment-energy revenue omitted from the compact notation is
\begin{align}
R_k^{\mathrm{deploy}}
={}&\sum_{t\in\calH_k}\lambda_t^{\RT}\dt
\left(\alpha_t^{\RU}c_t^{\RU,\actual}
-\alpha_t^{\RD}c_t^{\RD,\actual}\right.\nonumber\\
&\left.\hspace{18mm}+\alpha_t^{\SP}c_t^{\SP,\actual}
+\alpha_t^{\NSP}c_t^{\NSP,\actual}\right). \label{eq:app_deployment_revenue}
\end{align}

Equation~\eqref{eq:app_deployment_revenue} settles the expected signed AS energy at the RT LMP; kW times hours gives kWh and multiplication by dollars per kWh gives dollars.
No building/PV term appears in $R_k^{\mathrm{deploy}}$, DA/RT energy revenue,
or AS-capacity revenue. Their financial effect is already and exclusively
captured by the signed meter term and the demand epigraphs:
\begin{align}
C^{\mathrm{retail}}
={}&\sum_t C_t^{\TOU}
\bigl[p_t^{\EV}+p_t^{\BESS}
+\zeta^{\mathrm{BTM}}(p_t^{\load}-p_t^{\PV})\bigr]\dt\nonumber\\
&+C^{\PD}z^{\PD}+C^{\NCD}z^{\NCD}.
\label{eq:app_passive_retail_effect}
\end{align}

Equation~\eqref{eq:app_passive_retail_effect} calculates signed TOU cost and both incremental demand charges from the same complete meter. The nonlinear peak cost cannot generally be split into independent resource bills.
where the demand-charge increments are evaluated from the complete meter, not
from an independently settled passive account. Because those peak terms are
nonlinear and the EV--BESS optimizer can react to the exogenous profiles,
separate PV or building wholesale revenue would be misleading. All wholesale
energy, deployment, and capacity value is attributed only to EV and BESS.
For a full evaluation set $\calT$, the oracle problem is
\begin{align}
\min_{\boldsymbol{x}\in\mathcal{X}(\omega)}J_{\calT}
={}&\sum_{t\in\calT} C_t^{\TOU}p_t^{\GI}\dt
+C^{\PD}z^{\PD}+C^{\NCD}z^{\NCD}\nonumber\\
&-\sum_{t\in\calT}C^{\EV}p_t^{\EV}\dt\nonumber\\
&-\sum_{t\in\calT}\dt
\left(\lambda_t^{\DA}p_t^{\DA}
+\lambda_t^{\RT}p_t^{\RT}\right)\nonumber\\
&-\sum_{t\in\calT}\sum_{x\in\calK}\dt\,
\rho_t^{x,\DA}c_t^{x,\DA}\nonumber\\
&-\sum_{t\in\calT}\sum_{x\in\calK}\dt\,
\rho_t^{x,\RT}c_t^{x,\RT}
-R^{\mathrm{deploy}}. \label{eq:full_objective}
\end{align}

Equation~\eqref{eq:full_objective} collects retail energy and demand costs, subtracts EV-service receipts, subtracts DA/RT energy and capacity settlements, and subtracts expected deployment receipts. The displayed peak terms apply to one billing period; an evaluation spanning several months uses separate peak variables and rates for each month. The oracle sees the full evaluation data; MPC uses the corresponding truncated horizon and carried state.
The financial post-processor stores costs with negative signs. After independent
cent rounding by day, the validation accepts
\begin{align}
\left|V-\left(R^{\EV}+R^{\WM}+C^{\TOU}
+C^{\PD}+C^{\NCD}\right)\right|
&\le \$0.02. \label{eq:rounding}
\end{align}

Equation~\eqref{eq:rounding} is an ex-post accounting tolerance for independently rounded dollar amounts, not an optimization constraint or permission to omit a revenue component.
When passive inputs are enabled, the post-processor may expose their retail
cost contribution in NWM audit columns, but passive WM energy and capacity
accounts remain zero and the total financial identity is unchanged.

\section{Closed-loop MPC, billing states, and oracle dominance}
\label{app:closed_loop}

\subsection{Executed-state recursion}
Let $\xi_k$ collect the carried state before call $k$:
\begin{align}
\xi_k=\left(s_k,\boldsymbol{e}_k,M_k^{\PD,\hist},
M_k^{\NCD,\hist},\bar p_k^{\DA},\bar{\boldsymbol c}_k^{\DA}\right).
\label{eq:app_state_vector}
\end{align}
Given information $\mathcal{I}_k^\pi$, updater $\pi$ calls the truncated optimizer,
implements only $x_{k|k}^\pi$, and updates
\begin{align}
\xi_{k+1}=F_k(\xi_k,x_{k|k}^\pi,\omega_k). \label{eq:app_state_recursion}
\end{align}
The demand states in $F_k$ use
\begin{align}
M_{k+1}^{q,\hist}
=\max\{M_k^{q,\hist},m_q(k)p_k^{\GI,\mathrm{executed}}\}. \label{eq:app_threshold_recursion}
\end{align}
The DA components in \eqref{eq:app_state_vector} are submitted commitments and
therefore enter the corresponding RT call as constants.

\subsection{Why in-horizon Perfect information is not the upper bound}
Let $Q_H(\xi_k,a;\mathcal{I}_k)$ be the optimized $H$-step cost after fixing
first action $a$. The implemented 24-h controller selects
\begin{align}
a_k^\pi\in\arg\min_a Q_H(\xi_k,a;\mathcal{I}_k^\pi). \label{eq:app_truncated_action}
\end{align}
The full-period optimal action instead minimizes
\begin{align}
Q^\star(\xi_k,a;\omega)
=\ell_k(\xi_k,a;\omega_k)+J_{k+1}^\star(F_k(\xi_k,a,\omega_k);\omega).
\label{eq:app_true_q}
\end{align}
Here $\ell_k$ is not a new objective: it is the executed one-interval
contribution of \eqref{eq:objective}, i.e., signed TOU energy plus any
incremental PD/NCD charge minus EV charging-service, energy-market, AS-capacity,
and deployment-energy revenue. The term $J_{k+1}^\star$ is the exact minimum
cost from the next state to the end of the billing period. The 24-h controller
contains the former but not the latter.
Unless the truncated objective contains the exact continuation value,
$Q_H$ and $Q^\star$ need not order candidate actions identically. Perfect
information can improve $Q_H$ while choosing a first action that leaves a worse
SOC, demand threshold, remaining EV energy, or DA commitment outside the
horizon. The issue is therefore analogous to a greedy algorithm: local
information quality does not repair a missing global continuation value.

For the conditional dominance result, assume the full-period oracle and the
closed-loop policy share the same realized data, initial state, feasible set,
and terminal conditions.
The policy's complete executed trajectory is one feasible candidate in the
oracle problem, so
\begin{align}
J_{\mathrm{oracle}}^\star
&=\min_{\boldsymbol{x}\in\mathcal{X}(\omega)}J_{\calT}(\boldsymbol{x};\omega)
\le J_{\calT}(\boldsymbol{x}^{\pi};\omega), \label{eq:app_oracle_cost_proof}\\
V_{\mathrm{oracle}}^\star&\ge V^\pi,\qquad V=-J. \label{eq:app_oracle_value_proof}
\end{align}
This is a conditional upper-bound theorem, not an unconditional claim about
the weekly numerical benchmark, whose DA offer envelope can differ from
forecast-driven MPC. A finite-gap oracle incumbent is also distinct from its
certified objective bound.

\subsection{Demand-charge objective equivalence}
For fixed historical threshold $M$ and candidate horizon peak $z$,
\begin{align}
C^q\max\{M,z\}
=C^qM+C^q\max\{z-M,0\}. \label{eq:app_demand_equivalence}
\end{align}
The first term is constant within an optimization call, so the absolute and incremental
forms produce the same optimizer. The incremental form is used because executed
financial accounting should charge only a newly established billing peak. This
equivalence does not eliminate horizon truncation: a 24-h controller still does
not know which future interval will establish the final monthly maximum.

\section*{CRediT authorship contribution statement}
Yizhan Gu: Conceptualization, Methodology, Software, Data curation, Investigation, Validation, Visualization, Writing and editing. Wente Zeng: Conceptualization, Methodology, Resources, Project administration, Writing--review and editing. Hadi Eshraghi: Methodology, Supervision, Writing--review and editing. Jan Kleissl: Conceptualization,
Methodology, Supervision, Funding acquisition, Project administration, Writing--review and editing.

\section*{Funding}
This work was conducted within the UC San Diego--TotalEnergies research
collaboration. The authors retained responsibility for the study design, data
processing, analysis, interpretation, manuscript preparation, and decision to
submit the work for publication.

\section*{Data availability}
CAISO market-price and attenuation inputs and SDG\&E tariff documents are
available from the public sources cited in this article. Charger-session and
site-meter records are subject to the data providers' access restrictions;
this article does not imply permission to redistribute those records.

\section*{Declaration of competing interest}
Wente Zeng and Hadi Eshraghi are affiliated with TotalEnergies, which
participated in the research collaboration. The remaining authors declare no
known competing financial interests or personal relationships that could have
appeared to influence the work reported in this paper.

\section*{Declaration of generative AI and AI-assisted technologies in the manuscript preparation process}
During the preparation of this work, the authors used OpenAI ChatGPT and Codex
to support language editing, manuscript organization, code review, and
reproducible plotting from author-verified data. After using these tools, the
authors reviewed and edited the content as needed and take full responsibility
for the content of the published article. Generative-image tools were not used;
all manuscript figures were produced deterministically from the authors' data
and plotting code.

\section*{Acknowledgements}
The authors acknowledge the UC San Diego--TotalEnergies research collaboration
and the personnel who supported access to the charger-session, building-meter,
PV, tariff, and market data used in this study.

\clearpage
\bibliographystyle{elsarticle-num}
\bibliography{references_v3.3}

\end{document}
